% !TeX program = xelatex
\documentclass[UTF8,a4paper,oneside,openany,fontset=fandol,zihao=-4]{ctexbook}
\usepackage[top=22mm,bottom=22mm,left=22mm,right=22mm,headheight=16pt]{geometry}
\usepackage{amsmath,amssymb,amsthm,booktabs,longtable,array,enumitem}
\usepackage{tikz}
\usetikzlibrary{arrows.meta,positioning,calc}
\usepackage{xurl}
\usepackage[unicode,hidelinks]{hyperref}
\setmainfont{Latin Modern Roman}
\setsansfont{Latin Modern Sans}
\setmonofont{Latin Modern Mono}[Scale=MatchLowercase]
\hypersetup{pdftitle={Erdős 第650号问题：高中生可读的完整证明详解},pdfauthor={证明讲解与整理},pdfsubject={基于 arXiv:2603.28636v1 的逐步证明讲义}}
\setlength{\emergencystretch}{3em}
\setlist[enumerate]{leftmargin=2.2em,itemsep=.4em}
\setlist[itemize]{leftmargin=2em,itemsep=.3em}
\setcounter{tocdepth}{1}
\setcounter{secnumdepth}{2}
\allowdisplaybreaks[1]
\renewcommand{\proofname}{证明}
\newtheoremstyle{guide}{.7em}{.7em}{\normalfont}{}{\bfseries}{。}{.5em}{}
\theoremstyle{guide}
\newtheorem{theorem}{定理}[chapter]
\newtheorem{lemma}[theorem]{引理}
\newtheorem{proposition}[theorem]{命题}
\newtheorem{example}[theorem]{例}
\newtheorem{claim}[theorem]{断言}
\newtheorem{remark}[theorem]{说明}
\newcommand{\N}{\mathbb N}
\newcommand{\Z}{\mathbb Z}
\newcommand{\ceil}[1]{\left\lceil#1\right\rceil}
\newcommand{\floor}[1]{\left\lfloor#1\right\rfloor}
\DeclareMathOperator{\lcm}{lcm}
\begin{document}
\begin{titlepage}
\centering
\vspace*{18mm}
{\zihao{1}\bfseries Erd\H{o}s 第 650 号问题\par}
\vspace{10mm}
{\zihao{2}\bfseries 高中生可读的完整证明详解\par}
\vspace{13mm}
{\Large 从整除与配对，一步步走到最优答案\par}
\vspace{14mm}
\[
\boxed{\displaystyle f(m)=\min\left\{m,\left\lceil2\sqrt m\right\rceil\right\}}
\]
\vspace{16mm}
\begin{minipage}{.82\textwidth}
\large
\textbf{阅读起点：}会计算整除、最大公约数、平方与平方根，能看懂简单的集合与不等式。\par\medskip
\textbf{本书补齐：}配对准则的证明、开区间端点的处理、三类选点为什么不碰撞、中国剩余定理的用法与证明、精确取整的全部步骤。
\end{minipage}
\vfill
{原论文作者：Wouter van Doorn、Yanyang Li、Quanyu Tang\par}
\vspace{3mm}
{依据 arXiv:2603.28636v1；详解版整理日期：2026 年 11 月 23 日\par}
\end{titlepage}
\frontmatter
\chapter*{怎样读这份讲义}
\addcontentsline{toc}{chapter}{怎样读这份讲义}
这份讲义讲解的是 Erd\H{o}s 第650号问题，即怎样把一组正整数与某个开区间里的互不相同的倍数配起来。原论文为 Wouter van Doorn、Yanyang Li、Quanyu Tang 的 \emph{Optimal bounds for an Erd\H{o}s problem on matching integers to distinct multiples}\cite{paper}。核心定理和证明路线来自该论文；这里重新安排讲解顺序，补充所需的初等引理、逐行计算、反例、图示与练习。

“高中生可读”不表示每一页都容易。尤其是上界构造需要耐心处理余数和素因子。我们把这些知识在使用前证明出来，读者可以按顺序阅读，不必预先学过图论、Hall 定理或中国剩余定理。唯一默认采用的数论基础是整数的唯一质因数分解，以及初等整除运算。

推荐分三遍阅读。第一遍看第1章和每章开头，弄清楚问题与两条证明路线；第二遍读配对准则和下界，重点理解“用一对整数识别一个数”；第三遍读上界与取整，理解为什么保证值不能再提高。最后的具体例子和练习用于检查是否掌握了推理。

本讲义直接把开区间的端点例外纳入证明，而不把 $|\Gamma(S)|\ge2\sqrt{|S|}$ 错当作对原图所有情形都成立。原文中素数幂论证的严格不等号也在这里写明适用条件。本文是对既有数学证明的细化，不是另行发现一个新定理。随包的有限计算用于查错，不替代对所有整数的证明；本次没有重新运行外部 Lean 形式化工程。

\clearpage
\tableofcontents
\clearpage
\mainmatter
% ===== intro =====
\chapter{先把问题说清楚}
\section{给每个数挑一个自己的倍数}
从 $m$ 个\textbf{互不相同的正整数}开始，把它们从小到大排好：
\[
A=\{a_1<a_2<\cdots<a_m\},\qquad L=a_m.
\]
这里 $m$ 是数的个数，$L$ 是其中最大的数。例如 $A=\{2,5,9\}$ 时，$m=3$，$L=9$。它们承担不同的作用。

然后任选一个实数 $x$，只允许在开区间
\[
I=(x,x+2L)
\]
中挑整数。区间长度固定为 $2L$，两个端点都不能选。若 $x=0,L=4$，可用整数是 $1,2,\ldots,7$；$0$ 和 $8$ 都不在区间内。

我们希望找出一些数对 $(a,b)$，满足下面三个要求。
\begin{enumerate}
\item $a$ 来自 $A$，$b$ 是区间 $I$ 内的整数。
\item $a\mid b$，即存在整数 $k$ 使得 $b=ka$。
\item 不同数对的第一个数互不相同，第二个数也互不相同。
\end{enumerate}
例如 $(2,4)$ 与 $(4,4)$ 不能同时选，因为第二个数重复了。$(2,2)$ 与 $(4,4)$ 可以同时选。第一列和第二列是不同角色：允许一个数同时出现在两列；限制的是\emph{同一列}不能重复。

倍数中的系数 $k$ 可以是负数，也可以是零。例如 $3\mid-6$，每个正整数也都整除 $0$。因为 $x$ 可以是任意实数，后面的论证必须对负区间和跨过零的区间同样成立。

\section{一个小例子：随便贪心可能选错}
取 $A=\{2,4\}$，$I=(0,8)$。数字 $2$ 可以选择 $2,4,6$，数字 $4$ 只能选择 $4$。若先给 $2$ 分配 $4$，就把 $4$ 唯一能用的位置占用了；改成
\[
(2,2),\quad(4,4)
\]
便能配成两对。

\begin{center}
\begin{tikzpicture}[x=1cm,y=1cm,every node/.style={font=\small}]
\node at (0,2.3){待配的数};\node at (4,2.3){可用倍数};
\node[circle,draw,minimum size=7mm] (a) at (0,1.2){$2$};
\node[circle,draw,minimum size=7mm] (b) at (0,-.5){$4$};
\node[circle,draw,minimum size=7mm] (u) at (4,1.5){$2$};
\node[circle,draw,minimum size=7mm] (v) at (4,0){$4$};
\node[circle,draw,minimum size=7mm] (w) at (4,-1.5){$6$};
\draw[gray,dashed] (a)--(v);\draw[gray,dashed] (a)--(w);
\draw[very thick] (a)--(u);\draw[very thick] (b)--(v);
\node[align=center] at (2,-2.2){粗线是同时选择的两条配对；\\虚线是还可以选择的关系。};
\end{tikzpicture}
\end{center}

所以，“每个数都能找到倍数”还不够；还必须保证它们最后选中的倍数彼此不同。这正是配对准则要解决的问题。

\section{题目问的是最坏情形能保证多少对}
对于给定的 $A,x$，用 $F(A,x)$ 表示最多能配成多少对。题目要找的 $f(m)$ 是一个\textbf{统一保证值}：无论怎样选 $m$ 个正整数，无论怎样移动长度为两倍最大数的区间，都至少能配出 $f(m)$ 对；而把这个保证再提高 $1$ 就会遇到反例。

因此证明一个候选答案 $r$，必须把两件事都做完。
\begin{itemize}
\item \textbf{下界：}对每一个允许的 $A,x$，证明 $F(A,x)\ge r$。
\item \textbf{上界：}找出至少一个允许的 $A,x$，证明 $F(A,x)\le r$。
\end{itemize}
下界要求没有遗漏任何情形；上界只需要一个精心设计的情形。这里的上界是给“最坏情形的保证值”设限，并不是说每个具体问题都只能配这么多对。

\begin{example}
结论将给出 $f(6)=5$。但取 $A=\{1,2,3,4,5,6\}$、$I=(0,12)$ 时，可以选择 $(1,1),(2,2),\ldots,(6,6)$，共六对。这不与 $f(6)=5$ 冲突：六对是这个具体问题的最优值，五对才是面对所有可能问题都成立的保证。
\end{example}

\section{最终答案和符号}
\begin{theorem}[原论文的主结论]\label{thm:main}
对每个正整数 $m$，都有
\[
\boxed{f(m)=\min\{m,\ceil{2\sqrt m}\}.}
\]
\end{theorem}
其中 $\ceil y$ 表示不小于 $y$ 的最小整数，例如 $\ceil{4.2}=5$、$\ceil4=4$。另一个常用符号 $\floor y$ 表示不大于 $y$ 的最大整数；例如 $\floor{-1.2}=-2$。$|S|$ 表示有限集合 $S$ 的元素个数，而不是把集合取绝对值。

为什么答案里一定有 $m$？因为每个左边的数只能使用一次，任意配对都至多有 $m$ 对。当 $m\ge4$ 时，$2\sqrt m\le m$：两边非负，平方后等价于 $4m\le m^2$，再除以正数 $m$ 就是 $4\le m$。由于 $m$ 本身是整数，$\ceil{2\sqrt m}\le m$。所以 $m\ge4$ 时可以只写 $f(m)=\ceil{2\sqrt m}$；小的 $m$ 仍需保留这个最小值。

\section{证明的地图}
下界的核心是一种“编码”，但要先处理一个例外：如果 $x$ 是 $L$ 的整数倍，就先将最大数 $L$ 配给中点 $x+L$，暂时删去这两个顶点，最后再加回这一对。这样可以避开开区间端点造成的困难。

在一般情形的原图，或上述删点后的图里，给当前处理的每个 $a$ 选择区间左半边和右半边的两个倍数，把 $a$ 编码成一个有序数对。若不同的 $a$ 得到不同的数对，左边的可选整数有 $p$ 个、右边有 $q$ 个，就至多编码 $pq$ 个不同的 $a$。再用
\[
(p-q)^2\ge0
\quad\Longrightarrow\quad
pq\le\frac{(p+q)^2}{4},
\]
便能从乘积的下界得到总数 $p+q$ 的下界。配对准则负责把这个计数结果变成互不冲突的实际配对。

上界的核心是把 $st$ 个不同的数排成一个 $s\times t$ 数表，让每个数在指定区间内只能选择一个“行位置”或一个“列位置”。总共只有 $s+t$ 个位置，因此任何配对都不能超过 $s+t$ 对。余数构造负责把这张数表真正变成满足整除要求的整数。

\begin{center}
\begin{tikzpicture}[node distance=8mm,>=Latex,every node/.style={font=\small},box/.style={draw,rounded corners,align=center,text width=5.8cm,inner sep=6pt}]
\node[box] (l1) {一般情形直接编码\\端点对齐时先固定一对并删点};
\node[box,below=of l1] (l2) {$|S|\le pq\le(p+q)^2/4$\\得到邻域下界，再使用配对准则};
\node[box,below=of l2] (l3) {必要时加回固定的一对\\保证至少 $\min(m,\ceil{2\sqrt m})$ 对};
\node[box,right=12mm of l1] (r1) {构造特殊的 $s\times t$ 数表\\每个格子只对应一个行位置和一个列位置};
\node[box,below=of r1] (r2) {$st$ 个数，共用 $s+t$ 个位置\\得到 $f(st)\le s+t$};
\node[box,below=of r2] (r3) {选择最接近的整数参数\\得到相同的上界};
\draw[->] (l1)--(l2);\draw[->] (l2)--(l3);\draw[->] (r1)--(r2);\draw[->] (r2)--(r3);
\end{tikzpicture}
\end{center}
以上只是路线图。接下来的各章将证明每一条箭头，而不是把图中的结论当作已经成立。

% ===== matching =====
\chapter{怎样保证所有配对互不冲突}
\label{mat:chapter}

\section{一个人有选择，不等于所有人能同时选好}

先暂时把整除条件放在一边。设左边有一个有限集合 $A$，右边有一个有限集合 $B$；在允许配对的两个元素之间连线。这样的图称为\textbf{二部图}。两边是两组不同的对象；即使左右各有一个写着同样数字的点，也按所在的一侧区分。

一个\textbf{匹配}是一组不共用端点的连线。选择它，相当于同时完成若干配对，每个左点和每个右点最多参加一次。记 $n=|A|$，并用 $\nu(G)$ 表示最多可以同时完成的配对数。由于图是有限的，可选的连线集合只有有限多个，因而这个最大值一定取得。达到它的匹配称为\textbf{最大匹配}。

对任意 $S\subseteq A$，定义其\textbf{邻域}
\[
\Gamma(S)=\{b\in B:\text{至少有一个 }a\in S\text{ 可以与 }b\text{ 配对}\}.
\]
这里统计的是不同的右点，而非连线条数。如果三个左点都只能找同一个右点，那么每个左点都有选择，但三者合起来至多配成一对。因此，研究所有左点能否同时配好，必须考察每一个子集共有多少选择。

\begin{example}\label{mat:example-conflict}
设 $a_1$ 可以选择 $b_1,b_2$，$a_2$ 只能选择 $b_1$，而 $a_3,a_4$ 都只能选择 $b_3$。先配上 $a_1$ 与 $b_1$，并不意味着 $a_2$ 只能落单：把原配对改成 $a_1$ 与 $b_2$，就可以再配上 $a_2$ 与 $b_1$。不过，$a_3,a_4$ 中必有一个落单，因为它们共同只有 $b_3$ 这一个选择。于是最多能配成三对，且上述两对加上 $(a_3,b_3)$ 确实达到三对。
\end{example}

这个例子也区分了两件事：“在原有配对上不能直接再加一对”，未必意味着“配对总数已经最多”。有时需要先调整已有配对，才能多配一对。

\section{沿着一次改配看清交错路}

继续使用例~\ref{mat:example-conflict} 的连线。假设当前只选了 $(a_1,b_1)$ 和 $(a_3,b_3)$ 两对。$a_2$ 虽然尚未配好，却不能直接使用 $b_1$，因为它已经被 $a_1$ 占用。于是先顺着允许的连线从 $a_2$ 走到 $b_1$，再沿现有配对回到 $a_1$，最后走到尚未使用的 $b_2$。这条路依次经过
\[
a_2\longrightarrow b_1\longrightarrow a_1\longrightarrow b_2.
\]
第一条和第三条边是允许但尚未选用的边，中间一条是已经选用的边。把中间一对撤掉，改选两侧的两对，配对数便从两对增加到三对。沿路的内部点 $b_1,a_1$ 都只是换了伙伴；原来的落单点 $a_2$ 和空着的 $b_2$ 则各自得到一个伙伴。

换完后，只剩 $a_4$ 未配。它只能走到 $b_3$，再沿已有配对走到 $a_3$；但 $a_3$ 没有别的右侧选择，路线不能继续。能到达的左点是 $\{a_4,a_3\}$，能到达的右点是 $\{b_3\}$。这正好暴露了无法继续改配的原因：两个左点共用的候选只有一个。下面的证明把这种观察推广到任意有限二部图。

\section{配不上多少，恰好由最大的选择缺口决定}

定义集合 $S$ 的\textbf{亏损}
\[
d(S)=|S|-|\Gamma(S)|,
\qquad
D=\max_{S\subseteq A}d(S).
\]
亏损可以为负：这表示 $S$ 的右侧选择比左侧人数还多。但最大亏损 $D$ 不会为负，因为空集没有邻点，$d(\varnothing)=0$。保留空集，使“没有短缺”统一对应 $D=0$。

\begin{lemma}[匹配的亏损公式]\label{lem:hall-deficiency}
对任意有限二部图 $G$，若左点集 $A$ 有 $n$ 个元素，则
\begin{equation}\label{mat:deficiency-formula}
\nu(G)=n-\max_{S\subseteq A}\bigl(|S|-|\Gamma(S)|\bigr)=n-D.
\end{equation}
\end{lemma}

\begin{proof}
若 $n=0$，没有左点，匹配只能为空；唯一的子集也是空集。因此 $\nu(G)=0=D$，公式成立。以下考虑 $n\geq1$。

\textbf{第一步：任何配对都逃不开已有的亏损。}
固定一个匹配和一个子集 $S\subseteq A$。已经配上的 $S$ 中元素，必须使用 $\Gamma(S)$ 中互不相同的右点，因此至多有 $|\Gamma(S)|$ 个。于是，$S$ 中至少有 $|S|-|\Gamma(S)|$ 个左点没有配上；整个 $A$ 中没有配上的左点自然至少也有这么多。对所有 $S$ 取最大值，可知每个匹配至少留下 $D$ 个左点，所以
\begin{equation}\label{mat:easy-direction}
\nu(G)\leq n-D.
\end{equation}

\textbf{第二步：选定最大匹配，寻找可以重新安排配对的路线。}
取一个最大匹配 $M$，设未配上的左点组成集合 $U$。每条匹配边恰好配上一个左点，所以
\[
|U|=n-|M|=n-\nu(G).
\]
给所有连线规定方向：不属于 $M$ 的边，从左指向右；属于 $M$ 的边，从右指向左。从 $U$ 中任一点出发，沿这些箭头行走。到达的左点组成集合 $L$，到达的右点组成集合 $R$；起点本身也算到达，因此 $U\subseteq L$。若 $U$ 为空，则规定 $L=R=\varnothing$。

沿箭头走时，所经过的边交替为“未选进匹配的边”和“已选进匹配的边”。这种路线叫作\textbf{交错路}：先尝试把一个未配左点交给某个右点，再把该右点原来的左侧伙伴往后挪，继续尝试新的选择。

\textbf{第三步：每个可达右点都已经在 $M$ 中配好。}
假如 $R$ 中有一个没有配好的右点 $b$，那么从某个 $u\in U$ 到 $b$ 存在一条有向路。若路线重复经过一个点，就删去两次到达该点之间的绕行部分；这样可取得一条不重复顶点的路。

这条路从未配左点开始，在未配右点结束，第一条和最后一条边都不属于 $M$。设路上有 $r$ 条属于 $M$ 的边，则有 $r+1$ 条不属于 $M$ 的边。删去前者、加入后者，便得到一个新匹配：两个端点原本没有配对；每个内部顶点的原匹配边在路上，换边后仍然恰好参加一次配对；路外的配对不受影响。新匹配的边数为 $|M|-r+(r+1)=|M|+1$，与 $M$ 最大矛盾。因此，每个 $b\in R$ 都已配好。

\textbf{第四步：可达左点恰好由起点和这些右点的伙伴组成。}
对 $b\in R$，记它在 $M$ 中唯一的左侧伙伴为 $p(b)$。匹配边从 $b$ 指向 $p(b)$，所以 $p(b)\in L$，而且 $p(b)\notin U$。不同的右点有不同的匹配伙伴，因此 $p$ 是从 $R$ 到 $L\setminus U$ 的单射。

反过来，若 $a\in L\setminus U$，它不是起点，所以某条可达路线经过最后一条边进入 $a$。只有匹配边才从右侧指向左侧，故这条最后的边连接的右点正是 $a$ 的匹配伙伴，并且属于 $R$。因此 $p$ 也是满射。于是
\begin{equation}\label{mat:reachable-count}
|L|=|U|+|R|.
\end{equation}

\textbf{第五步：$L$ 的全部邻点恰好就是 $R$。}
每个 $b\in R$ 都通过某个左点走到，而那个左点属于 $L$，所以 $R\subseteq\Gamma(L)$。反过来，取 $a\in L$ 的任一邻点 $b$。若边 $ab$ 不属于 $M$，其方向是 $a\to b$，因此 $b\in R$。若边 $ab$ 属于 $M$，则 $a$ 不可能在 $U$ 中；根据第四步，$a$ 的匹配伙伴属于 $R$，而这个伙伴正是 $b$。所以也有 $\Gamma(L)\subseteq R$。两边合并得到
\[
\Gamma(L)=R.
\]
将它与式~\eqref{mat:reachable-count} 合用，就得到一个实际实现全部落单数的子集：
\[
d(L)=|L|-|\Gamma(L)|=|L|-|R|=|U|=n-\nu(G).
\]
由于 $D$ 是所有亏损的最大值，$D\geq d(L)=n-\nu(G)$，即 $\nu(G)\geq n-D$。再结合式~\eqref{mat:easy-direction}，两边相等，公式得证。若 $U=\varnothing$，这里的 $L,R$ 均为空，上述计数仍然成立。
\end{proof}

这个证明不仅算出最多配多少，还说明短缺出现在哪里：从所有未配左点沿交错路走到的左点集 $L$，就是一个亏损达到最大值的集合。当所有 $S$ 都满足 $|\Gamma(S)|\geq|S|$ 时，最大亏损为零，便能配好全部左点；这就是通常的 Hall 婚配定理。

也可以由证明得到一种寻找最大匹配的过程：从空匹配开始，只要找到从未配左点到未配右点的交错路，就沿路换边，每次多配一对。搜索时记录已到达的顶点，每个顶点只处理一次；因为顶点有限，每轮搜索都会找到所需路线，或在有限步内确认这样的路线不存在。总共只有 $n$ 个左点，成功增加的次数不可能超过 $n$。当这样的路已经不存在时，第三至第五步仍然成立，产生的集合 $L$ 说明任何匹配至多有当前这么多对。因此停止时确实得到最大匹配。这里只说明过程为何正确并会结束，不需要预先知道 Hall 定理。

\section{把平方根邻域界变成配对数}

\begin{proposition}\label{prop:neighbor-to-matching}
设有限二部图的左点集有 $n$ 个元素，$n$ 是非负整数。若每个非空子集 $S$ 都满足
\begin{equation}\label{mat:sqrt-neighborhood}
|\Gamma(S)|\geq2\sqrt{|S|},
\end{equation}
则
\[
\nu(G)\geq\min\bigl(n,\ceil{2\sqrt n}\bigr).
\]
\end{proposition}

\begin{proof}
当 $n=0$ 时，假设不涉及任何非空子集，结论为 $0\geq0$。以下设 $n\geq1$。

先证明整数序列 $h(t)=t-2\sqrt t$ 在 $t\geq1$ 时严格递增。对根式差作分母有理化，得到
\[
2\sqrt{t+1}-2\sqrt t
=\frac{2((t+1)-t)}{\sqrt{t+1}+\sqrt t}
=\frac{2}{\sqrt{t+1}+\sqrt t}<1.
\]
最后一步的理由是 $t\geq1$ 给出 $\sqrt{t+1}+\sqrt t\geq\sqrt2+1>2$。因此
\[
h(t+1)-h(t)=1-\frac{2}{\sqrt{t+1}+\sqrt t}>0.
\]
反复使用这个相邻项比较，便知对每个整数 $1\leq t\leq n$，都有 $h(t)\leq h(n)$。这里从 $1$ 开始比较；空集对应 $t=0$，须另行处理。

对任意非空 $S$，代入 $t=|S|$，由假设得到
\[
d(S)=|S|-|\Gamma(S)|
\leq |S|-2\sqrt{|S|}
\leq n-2\sqrt n.
\]
空集的亏损是零，故
\[
D\leq\max(0,n-2\sqrt n).
\]
利用引理~\ref{lem:hall-deficiency}，可得
\[
\nu(G)=n-D
\geq n-\max(0,n-2\sqrt n)
=\min(n,2\sqrt n).
\]
最后，$\nu(G)$ 为整数，必须不小于右边的向上取整值。又因 $n$ 是整数，
\[
\ceil{\min(n,2\sqrt n)}
=\min\bigl(n,\ceil{2\sqrt n}\bigr).
\]
这个等式可以分两种情况核对：若 $2\sqrt n\leq n$，则它的向上取整仍不超过整数 $n$；若 $2\sqrt n>n$，两边都等于 $n$。于是命题成立。
\end{proof}

小规模的情形也可以直接从亏损看清楚。当 $n=1,2,3$ 时，$n-2\sqrt n<0$；非空子集的亏损虽然为负，空集仍把最大亏损保持在零，因而全部左点都能配好。当 $n=5$ 时，最大亏损满足 $0\leq D\leq5-2\sqrt5<1$。由于 $D$ 是整数，只能有 $D=0$，仍然能配好五个左点。平方根给出的实数估计，正是借助这种整数限制变成了精确的配对保证。

\begin{remark}\label{mat:empty-set-warning}
空集的零亏损不能省略。例如只有一个左点、两个右点，且左点能选任一右点时，唯一非空子集的亏损是 $1-2=-1$。若错误地只取这个最大值，就会算出 $\nu=1-(-1)=2$；事实上一个左点最多配成一对。把空集也算进去，最大亏损为 $0$，正确答案就是 $1$。
\end{remark}

命题~\ref{prop:neighbor-to-matching} 给出的是一个充分条件。后面应用它时，需要对正在使用的那张图逐一验证非空子集的邻域界；如果删去了某些顶点，就应当使用删点后图中的邻域和左点数。这样，图论的计数结论才能准确接到整数的整除问题上。

\begin{example}\label{mat:why-all-subsets}
为什么不能只检查整个 $A$？设有九个左点和六个右点，其中前五个左点都只能选择 $b_1$，后四个左点都可以选择 $b_2,\ldots,b_6$。整个左点集的邻域有六个点，确实满足 $|\Gamma(A)|=6=2\sqrt9$。但前五个左点合起来最多配一对，后四个最多配四对，所以总共至多五对，达不到六对。问题就在子集 $S=\{a_1,\ldots,a_5\}$：它的邻域只有一个点，不满足命题的假设。检查每个子集，才能发现这种集中在局部的短缺。
\end{example}

\begin{example}\label{mat:sharp-graph-bound}
再看一个满足假设的例子：六个左点与五个右点之间的所有连线都允许。每个非空子集的邻域都是这五个右点；又因为 $1\leq|S|\leq6$，有 $2\sqrt{|S|}\leq2\sqrt6<5$，所以假设成立。直接把五个不同左点配给五个不同右点，可以配成五对；右点只有五个，无法配成六对。这与命题给出的 $\ceil{2\sqrt6}=5$ 完全一致。这里检验的是一般二部图的结论；后面的整除构造还会说明数论条件下如何达到相应的上界。
\end{example}

% ===== lower =====
\chapter{为什么任何情形都至少能配出这么多对}
\label{low:chapter}

上一章已经把匹配问题转化成邻域的计数问题。本章要完成剩下的关键一步：无论给出的正整数集合是什么，无论区间放在数轴的什么位置，都能找到足够多的匹配。这里的“无论”很重要。上界证明只需构造一个特殊例子；下界证明则必须覆盖所有允许的例子，尤其不能漏掉开区间的端点恰好是某个整数的倍数的情形。

阅读时可以始终追踪两个不同的问题：\textbf{选出的两个数是否合法，以及不同的原数能否得到同一个数对。}前一个问题靠区间不等式和整除关系解决；后一个问题靠比较坐标差，并在特殊位置进一步比较奇偶性解决。证明一个数对合法，并不自动证明全部数对互不相同。后文的朴素对称选点反例，将具体展示合法性检查通过而单射性检查失败的情况。

\section{先明确目标，以及计数与匹配的关系}

固定任意正整数 $m$，任取
\[
  A=\{a_1<a_2<\cdots<a_m\},\qquad L:=a_m,
\]
其中每个 $a_i$ 都是正整数。再任取实数 $x$，令
\[
  I=(x,x+2L),\qquad B=I\cap\Z.
\]
左边的点代表 $A$ 中的数，右边的点代表 $B$ 中的数；若 $a\mid b$，就在它们之间连边。这里 $a\mid b$ 的意思是：存在整数 $k$，使 $b=ka$。整数 $k$ 可以为正数、零或负数，因此 $0$ 与负整数也可能是合法的右边顶点。全文只要求除数 $a$ 为正，不要求被整除的数 $b$ 为正。

对于 $S\subseteq A$，它的邻域记作
\[
  \Gamma(S)=\{b\in B:\text{至少存在一个 }a\in S\text{ 满足 }a\mid b\}.
\]
所以 $|\Gamma(S)|$ 数的是不同的右边整数。同一个整数即使被 $S$ 中好几个数整除，在这里也只计一次。

这里也要区分“每个左点都有候选”和“不同左点能选到不同候选”。例如，两个左点各自都有一条边，但它们若都只连向同一个右点，就最多选出一对匹配。另一方面，只看所有左点合起来的候选总数也不够：可能有一小群左点争用少量候选，而其他左点拥有很多不能分给它们的候选。因此匹配定理要求检查每个子集；证明中“任意非空子集”这几个字，保证了局部的候选不足也会被纳入检查。

\begin{theorem}\label{low:main}
上述二部图中存在一个匹配，其大小至少为
\[
  \min\bigl(m,\ceil{2\sqrt m}\bigr).
\]
因此
\[
  f(m)\geq\min\bigl(m,\ceil{2\sqrt m}\bigr).
\]
\end{theorem}

根据命题~\ref{prop:neighbor-to-matching}，如果每个非空 $S\subseteq A$ 都满足
\begin{equation}\label{low:neighbor-target}
  |\Gamma(S)|\geq 2\sqrt{|S|},
\end{equation}
就能得到定理要求的匹配。这一命题来自引理~\ref{lem:hall-deficiency}：它把“任何一群左点共有多少个候选右点”转化成“整个图至少能选出多少条互不占用端点的边”。

本章的计数办法是：给每个 $a\in S$ 分配一个有序数对 $(v_a,w_a)$，两个数都是 $a$ 的倍数，而且分别位于区间中点的左右。随后证明不同的 $a$ 不会得到相同的数对。\textbf{这些数对本身不是最终的匹配。}不同数对仍然可以共用某一个坐标；它们的作用是提供邻域大小的下界，然后由前述匹配定理完成最后一步。

“有序”表示左右位置不能交换：第一坐标始终来自左半区，第二坐标始终来自右半区。比较两个数对是否相等，要求第一坐标相等且第二坐标也相等。正因为顺序已经固定，我们才能统一用“第二坐标减第一坐标”恢复有关原数的信息。

不过，式~\eqref{low:neighbor-target}不能在所有位置的原图上成立。我们会先处理一般位置，再单独处理端点对齐的特殊位置。这样分情况，是为了准确处理开区间不能包含端点的要求。

\section{用向下取整选出中点左边的最后一个倍数}

记区间中点为 $b=x+L$。对每个 $a\in A$，定义
\begin{equation}\label{low:floor-choice}
  u_a:=a\floor{\frac{x+L}{a}}=a\floor{\frac ba}.
\end{equation}
向下取整 $\floor t$ 是不超过实数 $t$ 的最大整数。若 $k=\floor t$，它的定义就是
\[
  k\leq t<k+1.
\]
取 $t=b/a$，再乘以正数 $a$，得到
\[
  u_a\leq b<u_a+a.
\]
把右侧的不等式移项，还得到 $b-a<u_a$。三者合起来是
\begin{equation}\label{low:floor-bounds}
  b-a<u_a\leq b<u_a+a.
\end{equation}

这组不等式同时说明了三个事实。第一，$u_a$ 是整数倍 $a\floor{b/a}$，因此确实是 $a$ 的倍数。第二，因为 $a\leq L$，所以
\[
  u_a>b-a=x+L-a\geq x,
\]
从而 $u_a$ 严格位于左端点右侧。第三，若另一个倍数 $ka$ 不超过 $b$，则 $k\leq b/a$；由于 $k$ 为整数，便有 $k\leq\floor{b/a}$，进而 $ka\leq u_a$。因此 $u_a$ 正是 $(x,b]$ 中最大的 $a$ 的倍数，既证明了它存在，也证明了“最大”这一性质。

注意这里使用的是向下取整，不是把小数部分直接删去。例如 $\floor{-5/4}=-2$，所以当 $b=-5$、$a=4$ 时，选到的是 $u_a=-8$，下一倍数是 $-4$，确实满足 $-8\leq-5<-4$。因此式~\eqref{low:floor-choice}在区间位于负半轴时仍然适用，没有暗中假定 $x$ 或 $b$ 为正。

从数轴上看，$a$ 的所有倍数排列成一列等距的点，相邻两点相距 $a$。我们选择中点左侧最后一个点，再选择它右侧紧邻的点。中点与左点的距离小于 $a$；这正是 $b-a<u_a$ 中严格不等号的来源。如果中点本身是倍数，就选择中点作为左点，而不是跳过它。先把这个规则固定，后面才有办法准确说明哪一种情形必须改动。

\section{一般位置：相邻倍数直接给出单射}

现在假设 $x\notin L\Z$，也就是不存在整数 $k$ 使 $x=kL$。把右边顶点分成两部分：
\[
  B_-=(x,b]\cap\Z,\qquad
  B_+=(b,x+2L)\cap\Z.
\]
中点若是整数，就只归入左半部分，所以两部分不相交，合起来恰好是 $B$。

式~\eqref{low:floor-bounds}已经说明 $u_a\in B_-$，并且 $u_a+a>b$。还需证明下一倍数没有落到右端点上。我们先有
\[
  u_a+a\leq b+a\leq b+L=x+2L.
\]
假如最后出现等号 $u_a+a=x+2L$，那么上述两步比较必须同时取等，即 $u_a=b$ 且 $a=L$。由于 $u_a$ 是 $a=L$ 的倍数，$b$ 就是 $L$ 的整数倍；又因为 $x=b-L$，$x$ 也会是 $L$ 的整数倍。这与本节假设矛盾。因此实际成立的是严格不等式
\[
  b<u_a+a<x+2L,
\]
所以 $u_a+a\in B_+$。

这段端点检查还说明一个细节：只要 $a<L$，就已经有 $u_a+a\leq b+a<b+L$，无须对 $x$ 作额外限制。真正可能触及右端点的是最大的那个数 $a=L$；而它只有在中点本身也是 $L$ 的倍数时才会触及端点。因此将特殊位置写成 $x\in L\Z$，并不是随意分类，而是由等号什么时候可能发生直接推出来的。

固定任意非空子集 $S\subseteq A$，记
\[
  \Gamma_-(S)=\Gamma(S)\cap B_-,\qquad
  \Gamma_+(S)=\Gamma(S)\cap B_+.
\]
对每个 $a\in S$，两个坐标 $u_a$ 与 $u_a+a$ 都是 $a$ 的倍数，因此定义
\[
  \phi:S\longrightarrow\Gamma_-(S)\times\Gamma_+(S),
  \qquad \phi(a)=(u_a,u_a+a)
\]
是合法的。若 $\phi(a)=\phi(a')$，两个有序数对的对应坐标相等，于是它们的第二坐标减第一坐标也相等，即
\[
  (u_a+a)-u_a=(u_{a'}+a')-u_{a'}.
\]
左边等于 $a$，右边等于 $a'$，故 $a=a'$。这就证明了 $\phi$ 是单射：数对中保留了能够唯一恢复原数 $a$ 的信息。

设 $p=|\Gamma_-(S)|$、$q=|\Gamma_+(S)|$。共有 $pq$ 个可能的有序数对，而 $|S|$ 个不同的数被送到不同数对，所以
\[
  |S|\leq pq.
\]
另一方面，$(p-q)^2\geq0$ 给出
\[
  (p+q)^2=p^2+2pq+q^2\geq4pq\geq4|S|.
\]
由于 $p+q$ 为非负数，开平方后得到
\[
  |\Gamma(S)|=p+q\geq2\sqrt{|S|}.
\]
$S$ 是任意选取的，因此命题~\ref{prop:neighbor-to-matching}适用，一般位置的证明到此完成。

这个计数过程没有要求我们选出的坐标用尽全部邻居。我们只是把每个 $a$ 放进由所有左邻居与所有右邻居组成的“数对集合”中；一个单射存在，就说明这个数对集合至少装得下 $|S|$ 个不同对象。至于它还有多少没有用到的数对，并不影响不等式。平方比较则说明：如果左右邻居总数很少，两者的乘积也不可能很大。因子 $2$ 就是在从乘积转回两部分之和时出现的。

\begin{example}\label{low:ordinary-example}
取 $A=\{3,6,10\}$、$L=10$、$x=1$，中点为 $b=11$，整个开区间是 $(1,21)$。三个数对应
\[
  \phi(3)=(9,12),\qquad
  \phi(6)=(6,12),\qquad
  \phi(10)=(10,20).
\]
它们的坐标差分别为 $3,6,10$，因此三个有序数对不同。但前两个数对的右坐标都是 $12$，所以不能把“每行都选择右坐标”当作匹配算法。

若取 $S=A$，完整的左右邻域分别为
\[
  \Gamma_-(S)=\{3,6,9,10\},\qquad
  \Gamma_+(S)=\{12,15,18,20\}.
\]
每边都有四个候选，而选点映射只使用其中一部分。这个例子中可以另行选出匹配 $(3,12)$、$(6,18)$、$(10,20)$；三个右点不同，每个右点又都是对应左点的倍数。一般证明不必事先找到这样的具体选法，因为匹配定理已经保证它存在。
\end{example}

\section{特殊位置：先把最大整数配给中点}

下面假设 $x\in L\Z$。设 $x=kL$，其中 $k\in\Z$，并记
\[
  b_0=x+L=(k+1)L.
\]
此时 $L$ 在开区间 $(x,x+2L)$ 中只有一个倍数，就是 $b_0$。事实上，若 $tL$ 位于这个区间内，除以正数 $L$ 后得到 $k<t<k+2$；整数 $t$ 只能等于 $k+1$。两旁的倍数 $kL$ 和 $(k+2)L$ 恰好是被排除的端点。

所以，若在原图中取 $S=\{L\}$，便有 $|\Gamma(S)|=1$，而 $2\sqrt{|S|}=2$。这明确说明：此时不能直接要求原图的每个子集都满足式~\eqref{low:neighbor-target}。解决办法是先保留合法配对
\[
  L\longleftrightarrow b_0,
\]
然后从左右两侧分别删去它们，得到
\[
  A_0=A\setminus\{L\},\qquad B_0=B\setminus\{b_0\}.
\]
剩下的图记为 $G_0$，只保留两端都还在的边。这样 $G_0$ 中的任何匹配都不会占用 $L$ 或 $b_0$，最后可以安全地加回刚才保留的这一对。

这里是在主动构造一类匹配，并没有声称每个最大匹配都必须先按这个顺序选边。为了证明“至少存在一个足够大的匹配”，找到一种能够保证成功的构造方式就够了。删除两个顶点后，剩余图中的边仍然保持原来的整除关系；我们既没有新增边，也没有让一个已经被占用的右点再次参加匹配。

若 $m=1$，集合 $A$ 只有 $L$，这一个配对已经达到目标
\[
  1=\min\bigl(1,\ceil{2\sqrt1}\bigr).
\]
因此以下只需讨论 $m\geq2$。此时 $A_0$ 有 $m-1$ 个元素，并且每个 $a\in A_0$ 都满足严格不等式 $0<a<L$。

还有一点不能混淆：删去 $L$ 后，本章中的字母 $L$ 仍表示原集合的最大数，区间也仍是原来那个长度为 $2L$ 的区间。我们不会把 $L$ 改成 $A_0$ 的最大数，也不会未经证明就套用关于 $f(m-1)$ 的结论。后面使用的是已经独立证明的图论命题，它只需要剩余图满足相应的邻域不等式。因此这一步不构成对下界结论的循环论证。

删去中点以后，左右两部分都成为开区间：
\[
  \widetilde B_-=(b_0-L,b_0)\cap\Z,\qquad
  \widetilde B_+=(b_0,b_0+L)\cap\Z.
\]
对 $S\subseteq A_0$，用 $\Gamma_0(S)$ 表示它在 $G_0$ 中的邻域，再记
\[
  \widetilde\Gamma_-(S)=\Gamma_0(S)\cap\widetilde B_-,\qquad
  \widetilde\Gamma_+(S)=\Gamma_0(S)\cap\widetilde B_+.
\]
我们仍准备把每个 $a$ 送到一个左右各取一数的有序数对。困难在于：如果 $a\mid b_0$，原来选出的 $u_a$ 就是已被删掉的中点，必须改选。

\section{为什么朴素的对称选点会失败}

一种自然的尝试是：当 $a\nmid b_0$ 时仍取相邻倍数 $(u_a,u_a+a)$；当 $a\mid b_0$ 时，改取围绕中点的 $(b_0-a,b_0+a)$。后者的两点确实都合格，但这还不够，因为不同的 $a$ 可能得到同一个数对。

\begin{example}\label{low:naive-collision}
取 $L=10$、$x=0$，于是 $b_0=10$，并考虑 $A=\{2,4,10\}$。删去 $10$ 后，$a=4$ 不整除中点，且
\[
  u_4=4\floor{10/4}=8,
  \qquad (u_4,u_4+4)=(8,12).
\]
而 $a'=2$ 整除中点，朴素对称方法给出的数对是
\[
  (b_0-2,b_0+2)=(8,12).
\]
不同的左点 $4$ 和 $2$ 得到了同一个数对，单射要求因此失败。问题并非某个坐标不合法：$8$ 和 $12$ 都在所需区间内，也同时是 $4$ 与 $2$ 的倍数。
\end{example}

这个例子也提示了修补方向。相邻倍数数对的坐标差为 $a$；由除数 $d$ 生成的对称数对的坐标差为 $2d$。两者相同时必有 $a=2d<L$。又因为相邻倍数的左坐标是 $a$ 的倍数，可写成 $2jd$，其中心就是
\[
  b_0=2jd+d=(2j+1)d.
\]
因此 $d\mid b_0$，但 $2d\nmid b_0$。这说明所有可能产生这类碰撞的除数，都落在同时满足
\[
  d\mid b_0,\qquad 2d<L,\qquad 2d\nmid b_0
\]
这三个条件的一组数中。对它们，把左右两点各向外再移动一个 $d$，改用 $(b_0-2d,b_0+2d)$。条件 $2d<L$ 保证移动之后仍未碰到开区间端点。

在例~\ref{low:naive-collision}中，$d=2$ 满足这些条件，改后的数对是 $(6,14)$，与 $a=4$ 对应的 $(8,12)$ 区分开了。不过，修好一个例子还不代表一般情况已被证明。下面必须同时检查：新的数对全部合法，而且不会制造其他类型的碰撞。

上面得到的三个条件，是这类潜在碰撞必须满足的条件，不是说每个满足条件的除数都一定已经发生碰撞。例如相应的 $2d$ 可能根本不在集合 $A_0$ 中。我们仍统一调整这些除数，是为了让定义只依赖 $d,L,b_0$，而不必再逐一询问某个竞争者是否存在。统一调整后的正确性，由接下来对所有类型组合的检查负责。

\section{三种类型的完整定义与落点检查}

对每个 $a\in A_0$，仍令 $u_a=a\floor{b_0/a}$，再按以下条件分类：
\begin{align*}
  \text{(T1)}\quad& a\nmid b_0;\\
  \text{(T2)}\quad& a\mid b_0,\quad 2a<L,\quad 2a\nmid b_0;\\
  \text{(T3)}\quad& a\mid b_0,\quad
       \bigl(2a\geq L\text{ 或 }2a\mid b_0\bigr).
\end{align*}
这是一个不重不漏的分类。首先，$a$ 要么整除 $b_0$，要么不整除。若整除，则“$2a<L$ 且 $2a\nmid b_0$”成立时归入 (T2)；该条件不成立，就表示其中至少一个要求不满足，即 $2a\geq L$ 或 $2a\mid b_0$，归入 (T3)。因此每个 $a$ 恰好属于一类。

定义
\begin{equation}\label{low:three-type-map}
  \phi(a)=
  \begin{cases}
    (u_a,u_a+a),&a\text{ 属于 (T1)},\\
    (b_0-2a,b_0+2a),&a\text{ 属于 (T2)},\\
    (b_0-a,b_0+a),&a\text{ 属于 (T3)}.
  \end{cases}
\end{equation}

\paragraph{(T1) 的两个点。}
由于 $a\nmid b_0$，$u_a$ 不会等于 $b_0$。将式~\eqref{low:floor-bounds}与 $a<L$ 合用，得到
\[
  b_0-L<b_0-a<u_a<b_0<u_a+a\leq b_0+a<b_0+L.
\]
两个坐标分别落在左右开半区，并且它们由相邻的整数倍组成，所以都是 $a$ 的倍数。

\paragraph{(T2) 的两个点。}
条件 $0<2a<L$ 给出
\[
  b_0-L<b_0-2a<b_0<b_0+2a<b_0+L.
\]
又因为 $a\mid b_0$，存在整数 $j$ 使 $b_0=ja$，于是 $b_0\pm2a=(j\pm2)a$，两个点都是 $a$ 的倍数。

\paragraph{(T3) 的两个点。}
此时 $0<a<L$，所以
\[
  b_0-L<b_0-a<b_0<b_0+a<b_0+L.
\]
同样将 $b_0$ 写成 $ja$，便有 $b_0\pm a=(j\pm1)a$，故这两个点也都是 $a$ 的倍数。

至此，对于任意 $S\subseteq A_0$ 及 $a\in S$，$\phi(a)$ 的左坐标属于 $\widetilde\Gamma_-(S)$，右坐标属于 $\widetilde\Gamma_+(S)$。我们检查的是严格不等式，因此没有把中点或两端点误收进来。

\begin{example}\label{low:type-table}
取 $L=12$、$x=0$、$b_0=12$，并令 $A=\{1,2,\ldots,12\}$。删去 $12$ 后，三种类型的实际取值如下。
\begin{center}
\renewcommand{\arraystretch}{1.12}
\begin{tabular}{c|c|c|c|c}
  $a$ & $u_a$ & 类型 & $\phi(a)$ & 右坐标减左坐标\\
  \hline
  $1$  & $12$ & (T3) & $(11,13)$ & $2$\\
  $2$  & $12$ & (T3) & $(10,14)$ & $4$\\
  $3$  & $12$ & (T3) & $(9,15)$  & $6$\\
  $4$  & $12$ & (T2) & $(4,20)$  & $16$\\
  $5$  & $10$ & (T1) & $(10,15)$ & $5$\\
  $6$  & $12$ & (T3) & $(6,18)$  & $12$\\
  $7$  & $7$  & (T1) & $(7,14)$  & $7$\\
  $8$  & $8$  & (T1) & $(8,16)$  & $8$\\
  $9$  & $9$  & (T1) & $(9,18)$  & $9$\\
  $10$ & $10$ & (T1) & $(10,20)$ & $10$\\
  $11$ & $11$ & (T1) & $(11,22)$ & $11$
\end{tabular}
\end{center}
例如 $a=4$ 属于 (T2)，因为 $4\mid12$、$8<12$，但 $8\nmid12$；$a=6$ 属于 (T3)，因为 $2a=12\geq L$。表内左坐标全在 $(0,12)$，右坐标全在 $(12,24)$。

同一张表还能再次看到修补的作用：如果 $a=4$ 仍使用朴素对称数对 $(8,16)$，它就会与 $a=8$ 的相邻倍数数对发生碰撞。实际定义把前者改为 $(4,20)$，后者保持为 $(8,16)$。选择向外移动时没有改变“坐标必须是原数的倍数”这一要求，因为移动的距离本身也是原数的整数倍。

表中允许某一坐标重复，例如 $a=2$ 与 $a=5$ 的左坐标都是 $10$；但两个完整数对 $(10,14)$ 与 $(10,15)$ 不同。这再次说明我们正在构造用于计数的单射，而不是直接构造匹配。数值表帮助理解定义，一般情形的单射性仍由下一节证明。
\end{example}

\section{逐一排除所有可能的数对碰撞}

\subsection{同一类型内不会碰撞}

对 (T1)、(T2)、(T3)，第二坐标减第一坐标依次为
\[
  a,\qquad 4a,\qquad 2a.
\]
若 $a$ 与 $a'$ 属于同一类型而 $\phi(a)=\phi(a')$，比较坐标差就分别得到 $a=a'$、$4a=4a'$ 或 $2a=2a'$。后两式除以相同的非零整数，也得到 $a=a'$。因此不同的同类元素不可能碰撞。

\subsection{不同类型之间，只需检查三种组合}

三种类型之间共有三种不计顺序的组合：(T1) 与 (T2)、(T1) 与 (T3)、(T2) 与 (T3)。交换两个元素的名称不会改变数对是否相等，所以检查这三种组合就覆盖了全部情况。

(T2) 和 (T3) 的两个坐标都关于 $b_0$ 对称，坐标和都等于 $2b_0$。因此，只要 (T1) 的元素 $a$ 与其中一类发生碰撞，就必须有
\begin{equation}\label{low:collision-midpoint}
  u_a+(u_a+a)=2b_0,
  \qquad\text{即}\qquad b_0=u_a+\frac a2.
\end{equation}
下面同时使用“坐标差”和“坐标平均值”这两条信息。

请注意，中点等式是\emph{假设碰撞之后}才得到的附加信息，并不是说每个 (T1) 数对原本都关于中点对称。例如前面的数对 $(9,12)$ 的平均值是 $10.5$，不等于当时的中点 $11$。只有它若与一个对称数对完全相同，两个坐标的平均值才被迫等于那个对称数对的中心。把这条附加信息和坐标差一起使用，才能产生下面的整除矛盾。

\paragraph{第一种组合：(T1) 与 (T2)。}
假设 $a$ 属于 (T1)，$a'$ 属于 (T2)，且两者数对相等。比较坐标差，得到 $a=4a'$。由于 $u_a$ 是 $a$ 的倍数，可写成 $u_a=ja=4ja'$，其中 $j\in\Z$。代入式~\eqref{low:collision-midpoint}，有
\[
  b_0=4ja'+\frac{4a'}2
     =4ja'+2a'
     =2a'(2j+1).
\]
$2j+1$ 是整数，因而这条等式说明 $2a'\mid b_0$。但 (T2) 的定义要求 $2a'\nmid b_0$，两者矛盾。

\paragraph{第二种组合：(T1) 与 (T3)。}
假设 $a$ 属于 (T1)，$a'$ 属于 (T3)，且两者数对相等。比较坐标差，得到 $a=2a'$。因为 $a\in A_0$，有 $a<L$，故 $2a'<L$。(T3) 定义中的“$2a'\geq L$ 或 $2a'\mid b_0$”必须成立；前一个条件已被排除，于是只能有 $2a'\mid b_0$。

另一方面，写 $u_a=ja=2ja'$，再利用相同的中点关系，得到
\[
  b_0=2ja'+\frac{2a'}2=(2j+1)a'.
\]
这说明 $b_0/a'$ 是奇整数。可是 $2a'\mid b_0$ 又要求存在整数 $h$，使 $b_0=2a'h$，即 $b_0/a'=2h$ 为偶整数。将两式相等并约去正数 $a'$，会得到 $2(h-j)=1$；左侧为偶整数，不可能等于 $1$，所以发生矛盾。

\paragraph{第三种组合：(T2) 与 (T3)。}
假设 $a$ 属于 (T2)，$a'$ 属于 (T3)，且两者数对相等。比较坐标差，得到
\[
  4a=2a',\qquad a'=2a.
\]
(T3) 的定义包含 $a'\mid b_0$，将 $a'=2a$ 代入便得到 $2a\mid b_0$。这与 $a$ 属于 (T2) 时要求的 $2a\nmid b_0$ 矛盾。

三种跨类组合全部排除，加上同类情形的结论，式~\eqref{low:three-type-map}定义的 $\phi$ 就是单射。

回看第二种组合，可以理解 (T3) 中两个条件各自的作用。如果 $2a'\geq L$，那么可能与它碰撞的 (T1) 元素必须等于 $2a'$，从而不可能属于所有元素都小于 $L$ 的集合 $A_0$；如果 $2a'<L$，就依靠 $2a'\mid b_0$ 排除碰撞。这里的“或”允许两个条件同时成立，并不要求只成立一个。把条件准确写成这样，才能在大小比较不够用时接着使用奇偶性比较。

\begin{remark}\label{low:zero-negative}
这些反证没有假定 $u_a$ 或 $b_0$ 为正。上面的 $j,h$ 都允许是负整数，奇偶性的定义“$2h$ 为偶整数、$2j+1$ 为奇整数”仍然成立。

还可以把中点改为负数，再按定义计算。取 $L=12$、$b_0=-12$，此时 $x=-24$。$a=4$ 仍属于 (T2)，因为 $4\mid-12$ 而 $8\nmid-12$，给出 $\phi(4)=(-20,-4)$。$a=5$ 属于 (T1)，并且
\[
  u_5=5\floor{-12/5}=5(-3)=-15,
  \qquad \phi(5)=(-15,-10).
\]
两组左坐标都在 $(-24,-12)$，右坐标都在 $(-12,0)$；整除关系并没有因为倍数为负数而失效。

若 $b_0=0$，任何正整数 $a$ 和 $2a$ 都整除 $0$，所以所有 $a\in A_0$ 都属于 (T3)，并且 $\phi(a)=(-a,a)$。由于 $a<L$，两个点分别在 $(-L,0)$ 和 $(0,L)$ 内；不同的 $a$ 给出不同数对。因此经过零点的特殊位置也已经包含在证明中。
\end{remark}

\section{由单射得到邻域下界，再把保留的一对加回来}

现在固定任意非空 $S\subseteq A_0$。已证明的单射给出
\[
  |S|\leq
  |\widetilde\Gamma_-(S)|\,
  |\widetilde\Gamma_+(S)|.
\]
设右侧两个因子分别为 $p$ 和 $q$。由于左右两部分不相交，并共同组成 $\Gamma_0(S)$，所以 $p+q=|\Gamma_0(S)|$。再次使用
\[
  (p+q)^2\geq4pq\geq4|S|,
\]
得到
\[
  |\Gamma_0(S)|\geq2\sqrt{|S|}.
\]
这次不等式作用于删去两个特殊顶点后的图 $G_0$，而不是原图。因为它对每个非空 $S\subseteq A_0$ 都成立，命题~\ref{prop:neighbor-to-matching}说明：$G_0$ 中有一个匹配，其大小 $r_0$ 满足
\[
  r_0\geq
  \min\bigl(m-1,\ceil{2\sqrt{m-1}}\bigr).
\]
这些计数全是有限集合的计数，因为一个有界区间中只能包含有限多个整数。空集也不会造成例外：它的邻域仍为空，所需不等式成为 $0\geq0$。使用图论命题时，我们已经先处理了 $m=1$，因此此处剩余左点数 $m-1$ 至少为一，满足前章命题的条件。
加回预留配对 $L\leftrightarrow b_0$，所得匹配没有端点冲突，而且大小增加 $1$。因此原图中有匹配大小 $r$ 满足
\begin{align}
  r&\geq1+\min\bigl(m-1,\ceil{2\sqrt{m-1}}\bigr)\notag\\
   &=\min\bigl(m,1+\ceil{2\sqrt{m-1}}\bigr).
   \label{low:add-one}
\end{align}
等号所用的事实是：给两个数都加上同一个数，较小者也恰好增加该数。

最后必须认真比较这里的 $m-1$ 与目标中的 $m$，不能把它们直接替换。对 $m\geq2$，利用平方差公式，有
\[
  2\sqrt m-2\sqrt{m-1}
  =\frac{2\bigl(m-(m-1)\bigr)}{\sqrt m+\sqrt{m-1}}
  =\frac2{\sqrt m+\sqrt{m-1}}<1.
\]
严格小于 $1$ 的理由是：分母至少为 $\sqrt2+1>2$。于是
\[
  2\sqrt m<1+2\sqrt{m-1}.
\]
向上取整保持大小关系，而且整数 $1$ 可以移到取整符号外，因此
\[
  \ceil{2\sqrt m}
  \leq\ceil{1+2\sqrt{m-1}}
  =1+\ceil{2\sqrt{m-1}}.
\]
具体来说，若 $u\leq v$，任何不小于 $v$ 的整数也不小于 $u$，所以 $\ceil u\leq\ceil v$；而整数集合整体平移 $1$，其最小合格整数也平移 $1$，所以 $\ceil{v+1}=\ceil v+1$。

把这条比较代入式~\eqref{low:add-one}，并注意与同一个 $m$ 取较小值不会颠倒不等号，便得到
\[
  r\geq\min\bigl(m,\ceil{2\sqrt m}\bigr).
\]
可以用小整数检查取整后的含义。当 $m=5$ 时，剩下的四个左点至少能配出四对，加回预留的一对得到五对，恰好等于目标 $\min(5,\ceil{2\sqrt5})=5$。当 $m=6$ 时，本特殊情形的证明实际上保证六对，而统一目标只有 $\min(6,\ceil{2\sqrt6})=5$。这并不矛盾：下界要求所有情形至少达到目标，某些位置能配出更多对完全允许。尤其不能因为某个特殊位置得到更强结果，就把它直接当成所有位置都成立的新结论。

$m=1$ 的特殊位置已经单独证明，一般位置也已完成；这覆盖了每个正整数 $m$、每个允许的集合 $A$ 和每个实数 $x$。因此无论怎样选择原问题的数据，都保证能配出至少这么多对，定理~\ref{low:main}得证。

% ===== upper =====
\chapter{怎样构造只能配出少量对的例子}\label{up:chapter}

\section{证明上界，需要找怎样的例子？}\label{up:goal}

回顾 $f(m)$ 的定义：无论怎样选择 $m$ 个不同的正整数组成集合 $A$，无论怎样选择长度为 $2\max A$ 的开区间，都能保证配出 $f(m)$ 对；每对中的右边整数是左边整数的倍数，而且左右两边各自不能重复使用。要证明一个保证值不能超过某个数，只需找出一个例子，使它配不出更多的对。这里的“一个例子”必须同时指定集合和区间，不能只指定其中之一。

\begin{theorem}\label{thm:product-upper}
对任意正整数 $s,t$，都有
\[
f(st)\leq s+t.
\]
\end{theorem}

本章把这个定理的证明拆开完成。若 $s=1$ 或 $t=1$，左边总共只有 $st$ 个整数，因此任何匹配至多有 $st$ 对，而 $st<s+t$。以下只需考虑 $s,t\geq2$。结论中的 $st$ 和 $s+t$ 在交换 $s,t$ 后都不改变，所以如果原来 $s>t$，交换二者的名字即可。我们据此假设
\[
2\leq s\leq t.
\]

我们准备把 $st$ 个整数排成 $s$ 行、$t$ 列，记为 $\alpha_{i,j}$。希望第 $i$ 行的所有整数共用一个倍数 $L_i$，第 $j$ 列的所有整数再共用一个倍数 $R_j$，而区间里除此之外没有其他倍数。这样，虽然左边有 $st$ 个不同整数，右边能够使用的整数却只有
\[
L_1,\ldots,L_s,\quad R_1,\ldots,R_t,
\]
一共 $s+t$ 个。每对必须占用一个不同的右边整数，于是匹配的对数至多为 $s+t$。注意，多个左边整数可以都整除同一个右边整数，但在一次匹配中，那个右边整数仍然只能用一次。

例如在三行三列的设想中，第一行的三个整数都可以把 $L_1$ 作为倍数，第二行的三个整数都可以把 $L_2$ 作为倍数，第三行同理。每一列还共同拥有自己的 $R_j$。乍看每个格子都有两个选择，九个格子合起来似乎有十八种选择，但这里计算的是选择的次数，其中许多选择指向同一个整数。匹配要求不同配对使用不同的目标，所以真正限制配对数量的是六个可用目标，而不是十八次出现。

把这个计数说得更严格：从任意一组合法配对中，将每一对送到它使用的右侧整数。由于配对规则禁止重复使用右侧整数，这个映射是单射。有限集之间存在单射时，出发集合的元素个数不超过到达集合的元素个数，因此配对数不超过右侧列表的总长度。这个论证不要求我们提前找出最大匹配，也不要求知道每一个右侧整数最终都会被用上。

这个安排的关键是让
\[
L_i=x_0-i,\qquad R_j=x_0+M+jD,\qquad
\alpha_{i,j}=M+i+jD.
\]
三者满足 $R_j-L_i=\alpha_{i,j}$。只要 $\alpha_{i,j}$ 整除 $L_i$，它就同时整除 $R_j$。因此，整个数论构造要解决的是：能否找到同一个 $x_0$，使所有要求
\[
x_0\equiv i\pmod{\alpha_{i,j}}
\]
同时成立？模数之间未必互素，所以在构造数字之前，我们先证明要用到的同余知识。

\section{同余工具：为什么许多余数要求可以同时满足？}\label{up:crt-tools}

\subsection{从整除到裴蜀等式}

符号 $d\mid n$ 的意思是存在整数 $k$ 使 $n=dk$；它同样适用于负整数 $n$。符号 $x\equiv r\pmod n$ 的意思是 $n\mid(x-r)$。这里 $r$ 不一定是从 $0$ 到 $n-1$ 的标准余数，用任意一个与它同余的整数代替，要求都不改变。两个正整数的最大公约数记为 $\gcd(a,b)$；当它等于 $1$ 时，称 $a,b$ 互素。

\begin{lemma}[裴蜀等式]\label{up:bezout}
对任意正整数 $a,b$，存在整数 $u,v$，使
\[
ua+vb=\gcd(a,b).
\]
\end{lemma}

\begin{proof}
考虑所有形如 $ua+vb$ 的正整数，其中 $u,v$ 可以取正数、零或负数。这样的数存在，例如取 $u=1,v=0$ 就得到 $a$。在这些正整数中选最小的一个，记为 $d=u_0a+v_0b$。

把 $a$ 除以 $d$，按带余除法写成 $a=kd+r$，其中 $0\leq r<d$。那么
\[
r=a-kd=(1-ku_0)a-kv_0b.
\]
如果 $r>0$，它就是比 $d$ 更小的、同样形式的正整数，与 $d$ 的选取矛盾。因此 $r=0$，所以 $d\mid a$。对 $b$ 做同样的带余除法，得到 $d\mid b$。另一方面，任何同时整除 $a,b$ 的正整数，都整除 $u_0a+v_0b=d$，因而不可能大于 $d$。故 $d$ 正是最大公约数，这就得到所需的等式。
\end{proof}

裴蜀等式也说明：如果 $a,b$ 互素，乘以 $a$ 不会使模 $b$ 的信息丢失。确切地说，若 $b\mid ac$，取 $ua+vb=1$ 后乘以 $c$，便有
\[
c=uac+vbc.
\]
右边两项都能被 $b$ 整除，所以 $b\mid c$。这里允许约去 $a$ 的原因是互素条件；例如 $6\mid2\cdot3$，却不能由此说 $6\mid3$。

接下来还要使用素数分解。其要点可以从上述结论得到：若素数 $p$ 整除 $ab$ 而不整除 $a$，由于 $p$ 的正约数只有 $1,p$，此时 $\gcd(p,a)=1$，所以 $p\mid b$。因此，一个素数整除若干整数的乘积，必整除其中至少一个因子。任何大于 $1$ 的整数，若不是素数，就拆成两个更小的正因数；继续拆分会停止，得到素数的乘积。若一个整数有两种这样的分解，第一种分解的某个素因子必与第二种分解的某个素因子相同，约去这两个因子后重复，最终各个素数及其出现次数都相同。这就说明了素数分解的唯一性。

对于非零整数 $n$，记 $v_p(n)$ 为 $|n|$ 的素数分解中 $p$ 的个数。例如 $72=2^3\cdot3^2$，所以 $v_2(72)=3$、$v_3(72)=2$、$v_5(72)=0$。由唯一分解，正整数 $d$ 整除非零整数 $n$，当且仅当对每个素数 $p$ 都有
\[
v_p(d)\leq v_p(n).
\]
于是，判断一个整除关系，可以逐个比较素数的指数。这也是后面最大公约数证明采用的方法。

\subsection{模数两两互素时的中国剩余定理}

\begin{lemma}[互素模数的中国剩余定理]\label{up:coprime-crt}
设正整数 $n_1,\ldots,n_h$ 两两互素，$r_1,\ldots,r_h$ 为任意整数。则存在整数 $z$，同时满足
\[
z\equiv r_i\pmod{n_i}\qquad(1\leq i\leq h).
\]
而且可以把满足这些要求的 $z$ 取得任意大。
\end{lemma}

\begin{proof}
模数为 $1$ 的条件对任何整数都成立，可以先删去。如果所有条件都被删去，任取整数即可。否则，设
\[
Q=n_1\cdots n_h,\qquad Q_i=Q/n_i.
\]
因为各模数两两互素，$Q_i$ 与 $n_i$ 互素：如果有素数同时整除二者，它整除 $Q_i$ 就必整除某个 $n_j$，其中 $j\ne i$，这与 $n_i,n_j$ 互素冲突。根据引理~\ref{up:bezout}，可以找到整数 $c_i,d_i$，使
\[
c_iQ_i+d_in_i=1.
\]
令
\[
z=\sum_{i=1}^h r_i c_iQ_i.
\]
固定一个下标 $k$。当 $i\ne k$ 时，$Q_i$ 含有因子 $n_k$，所以相应求和项模 $n_k$ 等于 $0$；当 $i=k$ 时，$c_kQ_k$ 模 $n_k$ 等于 $1$。因此整个和模 $n_k$ 等于 $r_k$。这对每个 $k$ 都成立，故所有同余条件同时满足。最后，对任意整数 $w$，$z+wQ$ 仍然满足全部条件；由于 $Q>0$，令 $w$ 充分大即可把解取得任意大。
\end{proof}

这里的公式把每个余数要求单独放在一个求和项里。第 $i$ 项在模 $n_i$ 时负责留下 $r_i$，在其他模数下则全部变成零。互素性保证我们能够构造出这种“只在一个模数下留下作用”的项，而不是仅仅凭空断言解存在。

例如，要求 $z$ 除以 $3$ 余 $2$、除以 $5$ 余 $3$。此时 $Q=15$，两个 $Q_i$ 分别为 $5,3$。因为 $2\cdot5$ 模 $3$ 余 $1$，而 $2\cdot3$ 模 $5$ 余 $1$，可以取两个系数都为 $2$，从公式得到
\[
z=2\cdot2\cdot5+3\cdot2\cdot3=38.
\]
直接检查，$38=3\cdot12+2=5\cdot7+3$。减去 $2Q=30$ 后得到较小的解 $8$；继续加减 $15$ 的整数倍，则得到更多解。后面的构造关心的是某个解存在并且可以足够大，不要求它一定是最小正解。

\subsection{模数不互素时：兼容条件为什么足够？}

如果模数有公因子，同余要求可能互相冲突。例如，一个整数不可能同时满足“模 $2$ 余 $0$”和“模 $4$ 余 $1$”，因为前一个要求它是偶数，后一个要求它是奇数。反过来，“模 $6$ 余 $1$”与“模 $10$ 余 $3$”在共同的因子 $2$ 上都要求奇数，并不冲突；整数 $13$ 就同时满足它们。下面的定理把这一判断推广到有限多个模数。

\begin{lemma}[一般模数的中国剩余定理]\label{up:general-crt}
给定正整数 $n_1,\ldots,n_h$ 和整数 $r_1,\ldots,r_h$，同余方程组
\[
z\equiv r_i\pmod{n_i}\qquad(1\leq i\leq h)
\]
有整数解，当且仅当对任意 $i,j$ 都有
\[
\gcd(n_i,n_j)\mid(r_i-r_j).
\]
\end{lemma}

\begin{proof}
先证明条件必要。如果 $z$ 是解，$\gcd(n_i,n_j)$ 同时整除 $z-r_i$ 和 $z-r_j$，于是也整除二者之差 $r_j-r_i$。整除一个整数与整除它的相反数等价，故得到定理中的条件。

再证明条件充分。只需考虑在某个 $n_i$ 的素数分解中出现的素数，它们只有有限多个。对于每个这样的素数 $p$，令
\[
e_p=\max_{1\leq i\leq h}v_p(n_i),
\]
并挑选一个达到最大值的下标 $j(p)$。暂时不直接处理原来的所有模数，而是对每个素数只保留一个最强的要求：
\[
z\equiv r_{j(p)}\pmod{p^{e_p}}.
\]
不同素数的幂两两互素，所以引理~\ref{up:coprime-crt} 给出一个同时满足这些新要求的整数 $z$。

现在检查任意一个原条件。固定 $i$，设 $p$ 在 $n_i$ 中的指数为 $e_i=v_p(n_i)>0$。由于 $e_i\leq e_p$，$p^{e_i}$ 同时整除 $n_i$ 与 $n_{j(p)}$，因此它整除二者的最大公约数。假设中的兼容条件于是给出
\[
p^{e_i}\mid\bigl(r_i-r_{j(p)}\bigr).
\]
另一方面，新条件给出 $p^{e_p}\mid(z-r_{j(p)})$，因而也有 $p^{e_i}\mid(z-r_{j(p)})$。把这两个整除关系相减，得到 $p^{e_i}\mid(z-r_i)$。这对 $n_i$ 中出现的每一个素数都成立；比较素数分解便得到 $n_i\mid(z-r_i)$，即原来的第 $i$ 个同余要求。若 $n_i=1$，此要求自动成立；如果全部模数都是 $1$，一开始直接取 $z=0$ 即可。因此所有情况都已覆盖。
\end{proof}

这份证明还解释了“两两兼容”为什么已经足够。对每一个素数，最高的幂决定了最详细的余数；较低的幂只检查这个余数是否与自己的要求一致。兼容条件完成了这项检查，剩下的不同素数之间再由互素的中国剩余定理组合。我们稍后只需证明一条最大公约数整除关系，就可以一次性满足 $st$ 个同余条件。

为了看清“保留最高幂”这一步，考虑三个要求
\[
z\equiv3\pmod8,\qquad z\equiv7\pmod{12},\qquad z\equiv7\pmod{18}.
\]
三个模数的分解为 $8=2^3$、$12=2^2\cdot3$、$18=2\cdot3^2$。对素数 $2$，最高幂是 $8$；对素数 $3$，最高幂是 $9$。因此只保留
\[
z\equiv3\pmod8,\qquad z\equiv7\pmod9.
\]
令 $z=3+8k$，第二个条件变成 $3-k\equiv7\pmod9$，可取 $k=5$，得到 $z=43$。现在检查被合并的模 $12$ 条件：$43=12\cdot3+7$，确实满足。为什么它没有被丢掉？因为余数 $3$ 与 $7$ 在共同的因子 $4$ 上相同，而两个余数 $7$ 在共同的因子 $3$ 上相同；这两部分合起来就是模 $12$ 的要求。这个例子展示了证明中先分解、再合并的实际含义。

\section{选出合适的整数组成数表}\label{up:construction}

\subsection{用最小公倍数统一处理小素数}

定义
\[
D=\lcm(1,2,\ldots,st).
\]
最小公倍数是能被这些数全部整除的最小正整数。这样的正整数存在，因为它们的乘积就是一个公倍数。我们将反复使用两件事：一是每个从 $1$ 到 $st$ 的整数都整除 $D$；二是 $D\geq st\geq2s>s$。后一结论来自 $st\mid D$，以及当前的假设 $t\geq2$。

我们的目标数字形如 $M+i+jD$。比较两个这样的数字，参数 $M$ 会消去，留下
\[
(M+i+jD)-(M+k+lD)=(i-k)+(j-l)D.
\]
若 $i\ne k$，把 $i-k$ 记作 $q$、$j-l$ 记作 $r$，则 $0<|q|<s$、$|r|<t$。因此，需要提前关注这些可能的差值 $q+rD$ 的素因子。

\subsection{有限个坏素数，以及怎样避开它们}

定义
\[
\mathcal P=\left\{p>st:\ \begin{array}{l}
p\text{ 是素数，且对某些整数 }q,r,\\
0<|q|<s,\ |r|<t,\ p\mid(q+rD)
\end{array}\right\}.
\]
称它们为坏素数，只是为了提醒我们：这些大素数可能妨碍后面的兼容条件，需要在选 $M$ 时避开。这个名字不附加其他数学含义。

先确认 $\mathcal P$ 有限。满足范围限制的整数对 $(q,r)$ 只有有限多个，而且每个 $q+rD$ 都非零。后一点不能省略：若 $r=0$，它等于非零的 $q$；若 $r\ne0$，则 $|rD|\geq D>|q|$，故
\[
|q+rD|\geq |rD|-|q|>0.
\]
一个非零整数只有有限多个素因子；将有限多个这样的素因子集合合并，仍然是有限集，所以 $\mathcal P$ 有限。若某个差值是 $1$ 或 $-1$，它没有素因子，对这个集合没有贡献。

固定 $p\in\mathcal P$。为了不让 $p$ 整除任何 $M+i+jD$，应禁止 $M$ 取以下剩余类：
\[
-i-jD\pmod p,\qquad 1\leq i\leq s,\quad 1\leq j\leq t.
\]
指标对一共 $st$ 个，因此被禁止的不同剩余类至多有 $st$ 个；有些指标可能给出相同余数，禁止的个数还会更少。模 $p$ 总共有 $p$ 个剩余类，而 $p>st$，所以至少有一个没有被禁止的余数。对每个坏素数分别选一个这样的余数，记作 $c_p$。

不同坏素数彼此互素。由引理~\ref{up:coprime-crt}，可以同时满足
\[
M\equiv c_p\pmod p\qquad(p\in\mathcal P).
\]
而且可以把 $M$ 取得足够大，使
\begin{equation}\label{up:M-size}
M>2s+2tD.
\end{equation}
具体地，若一个解为 $M_0$，则给它加上各坏素数之积的足够大的整数倍，所有余数都不变而大小不断增加。如果 $\mathcal P$ 是空集，就没有余数需要避开，直接选择满足~\eqref{up:M-size} 的整数 $M$。至此，我们同时获得了大小条件和所有避开条件。

只把 $M$ 取得很大，不能代替余数限制。比如在后面将计算的 $s=t=2,D=12$ 的情形，若随便取 $M=104$，虽然它超过了所需的大小界，却会出现
\[
\alpha_{1,1}=117,\qquad\alpha_{2,2}=130,
\qquad\gcd(117,130)=13.
\]
这两个格子分别要求 $x_0$ 模 $117$ 余 $1$、模 $130$ 余 $2$。前一个要求 $x_0$ 模 $13$ 余 $1$，后一个要求它模 $13$ 余 $2$，两者冲突，根本不存在共同的 $x_0$。因此，“足够大”和“避开某些余数”是两个不同的任务，我们用中国剩余定理先实现后一个任务，再保持余数不变地扩大数值。

\subsection{数表中的数字确实互不相同}

现在定义
\[
\alpha_{i,j}=M+i+jD,\qquad
A=\{\alpha_{i,j}:1\leq i\leq s,\ 1\leq j\leq t\}.
\]
所有数都是正整数。若 $\alpha_{i,j}=\alpha_{k,l}$，则
\[
i-k=(l-j)D.
\]
左边的绝对值小于 $s$，而 $s<D$；右边是 $D$ 的整数倍。唯一绝对值小于 $D$ 的 $D$ 的整数倍是 $0$，所以 $i=k$。代回等式后得到 $(l-j)D=0$，因为 $D>0$，故 $j=l$。因此不同格子放的是不同整数，$|A|=st$。同时，$\alpha_{i,j}$ 随两个指标分别严格增加，所以
\begin{equation}\label{up:maxA}
\max A=\alpha_{s,t}=M+s+tD.
\end{equation}

\section{关键兼容性：最大公约数为什么整除行号之差？}\label{up:gcd-section}

\begin{proposition}\label{up:gcd-compatibility}
对任意两组合法指标 $(i,j)$ 与 $(k,l)$，都有
\[
\gcd(\alpha_{i,j},\alpha_{k,l})\mid(i-k).
\]
\end{proposition}

\begin{proof}
若 $i=k$，右边为 $0$，每个正整数都整除 $0$，结论成立。以下设 $i\ne k$，并记
\[
q=i-k\ne0,\qquad r=j-l,\qquad
g=\gcd(\alpha_{i,j},\alpha_{k,l}).
\]
因为 $g$ 整除两个数，也整除它们的差，所以
\begin{equation}\label{up:g-divides-difference}
g\mid(q+rD).
\end{equation}
我们的任务是证明：$g$ 的每个素因子，在 $g$ 中出现的次数都不超过它在 $q$ 中出现的次数。

先排除大素因子。任取素数 $p\mid g$，则 $p\mid(q+rD)$，并且 $p\mid\alpha_{i,j}$。如果 $p>st$，由于 $q,r$ 正好满足坏素数定义中的范围限制，就有 $p\in\mathcal P$。但我们选择 $M$ 时已保证 $M\not\equiv-i-jD\pmod p$，即 $p\nmid\alpha_{i,j}$，这产生矛盾。因此，$g$ 的每个素因子都满足 $p\leq st$。

接着固定这样的素因子 $p$，记 $e=v_p(q)$。我们要证明 $D$ 至少含有 $e+1$ 个因子 $p$，即 $v_p(D)>e$。如果 $e=0$，由于 $p\leq st$，整数 $p$ 本身就在定义 $D$ 的列表里，所以 $p\mid D$，从而 $v_p(D)\geq1>0=e$。

如果 $e\geq1$，则由于 $0<|q|<s$ 和 $s\leq t$，有
\[
1\leq q^2<s^2\leq st.
\]
因此，正整数 $q^2$ 本身也是列表 $1,2,\ldots,st$ 中的一项。按照最小公倍数的定义，它整除 $D$，也就是 $q^2\mid D$。这里并不是“一个数较小就整除较大的数”，而是我们特意让 $D$ 被整个列表中的每个整数整除。于是
\[
v_p(D)\geq v_p(q^2)=2e>e.
\]
现在可以看出，开头选择 $s\leq t$ 有一个具体用途：它保证行号之差的平方仍然落在定义 $D$ 的整数范围内。列号之差 $r$ 不需要同样的平方估计，因为它始终乘着整个 $D$，已经继承了 $D$ 含有的素数幂。两个指标在构造中的作用因而有所不同，不能把这里关于 $q$ 的论证随意改成只比较 $|q|$ 与 $D$ 的大小。
最后的严格不等式依据是当前分支的 $e\geq1$：它给出 $2e-e=e>0$。仅有 $q\ne0$ 并不能保证 $e>0$，因此必须把 $e=0$ 与 $e\geq1$ 分开处理。

这里需要指数严格大于 $e$，而不仅仅是大于等于。原因是：若两项含有同样多的因子 $p$，相加后可能出现更多的因子 $p$。例如 $12=2^2\cdot3$ 和 $20=2^2\cdot5$ 都恰含两个因子 $2$，但 $12+20=32$ 含有五个因子 $2$。我们要确保加上 $rD$ 后，$q$ 原有的最低一层非零余数不会被抵消，所以必须让 $rD$ 至少多含一个因子 $p$。

两种情况下都得到 $p^{e+1}\mid D$。把 $q$ 写成 $p^eu$，其中 $p\nmid u$；又因为 $p^{e+1}\mid rD$，可把 $rD$ 写成 $p^{e+1}v$。即使 $r=0$，取 $v=0$ 也适用。于是
\[
q+rD=p^e(u+pv).
\]
括号里的整数不能被 $p$ 整除，否则减去 $pv$ 后会得到 $p\mid u$，与 $u$ 的选取矛盾。因此
\[
v_p(q+rD)=e=v_p(q).
\]
结合~\eqref{up:g-divides-difference}，得到 $v_p(g)\leq v_p(q)$。这对 $g$ 的每个素因子都成立；对于不整除 $g$ 的素数，左边的指数为零，也不大于右边。比较素数分解，就得到 $g\mid q=i-k$。若 $g=1$，没有素因子需要检查，且 $1\mid q$ 同样成立。
\end{proof}

这个论证中两个选择各负其责：$D$ 含有足够多的小素数幂，使小素数在差值中的指数受控；$M$ 则避开所有可能干扰的大素数。我们不是要求不同的 $\alpha_{i,j}$ 两两互素。同行数字完全可以有公因子，因为它们要求的余数相同；真正需要的是公因子不能与两个行号之间的余数要求冲突。

可以把最后的指数计算单独放在小数字上理解。取 $q=12,D=360,p=2$，则 $q$ 恰含两个因子 $2$，而 $D$ 含三个。无论整数 $r$ 是正、零还是负，都有
\[
12+360r=4(3+90r).
\]
括号里是奇数，所以和仍恰含两个因子 $2$。一般证明中的 $p^e(u+pv)$ 做的正是这件事：先提出共有的 $p^e$，再检查剩下的一项模 $p$ 不为零，而另一项模 $p$ 为零。这样既说明指数为什么不增加，也避免了对负数或零倍数作额外猜测。

现在将引理~\ref{up:general-crt} 用于所有方程
\begin{equation}\label{up:x0-congruences}
x_0\equiv i\pmod{\alpha_{i,j}}.
\end{equation}
任意两个方程的余数差正是 $i-k$，而命题~\ref{up:gcd-compatibility} 已证明其模数的最大公约数整除这个差。因此存在同时满足全部方程的整数 $x_0$。如果希望 $x_0$ 为正，可以给一个解加上所有模数的最小公倍数的足够大的整数倍；不过定理允许区间出现在实数轴的任何位置，正性不是证明所必需的。

\section{把全部倍数压到两个有限列表里}\label{up:interval}

令
\[
x=x_0-M,\qquad I=(x,x+2\max A),
\]
并使用本章开头提出的两个列表
\[
L_i=x_0-i\quad(1\leq i\leq s),\qquad
R_j=x_0+M+jD\quad(1\leq j\leq t).
\]
对于固定的格子 $(i,j)$，方程~\eqref{up:x0-congruences} 告诉我们 $\alpha_{i,j}\mid L_i$。又因为
\[
R_j-L_i=M+i+jD=\alpha_{i,j},
\]
所以 $R_j$ 也是 $\alpha_{i,j}$ 的倍数，而且这两个倍数是相邻的倍数。

先逐项核对它们确实位于开区间内。利用~\eqref{up:maxA}，左侧列表中的点满足
\[
L_i-x=M-i>0,
\]
以及
\[
(x+2\max A)-L_i=M+2s+2tD+i>0.
\]
第一个严格不等式来自 $i\leq s$ 和 $M>2s+2tD>s$；第二个式子的每一项都为正。右侧列表中的点满足
\[
R_j-x=2M+jD>0,
\]
以及
\[
(x+2\max A)-R_j=2s+(2t-j)D>0,
\]
因为 $j\leq t$，所以 $2t-j\geq t>0$。这样就排除了“候选点恰好落在开区间端点”的可能。

再核对前后相邻的其他倍数都在区间外。前一个倍数为 $L_i-\alpha_{i,j}$，并且
\[
(L_i-\alpha_{i,j})-x=-2i-jD<0.
\]
后一个倍数为 $R_j+\alpha_{i,j}$，并且
\begin{align*}
(R_j+\alpha_{i,j})-(x+2\max A)
&=M+i+2jD-2s-2tD\\
&=(M-2s-2tD)+i+2jD>0.
\end{align*}
最后一步使用~\eqref{up:M-size}。由于 $\alpha_{i,j}>0$，所有更早的倍数比 $L_i-\alpha_{i,j}$ 还小，所有更后的倍数比 $R_j+\alpha_{i,j}$ 还大。因此，$I$ 内恰有两个 $\alpha_{i,j}$ 的倍数，就是 $L_i,R_j$。

这对每个格子都成立，所以 $I$ 中所有能够作为 $A$ 中某个元素的倍数的整数，恰好组成
\[
\{L_1,\ldots,L_s\}\cup\{R_1,\ldots,R_t\}.
\]
其中的 $L_i$ 彼此不同，$R_j$ 也彼此不同；此外，所有 $L_i<x_0$，所有 $R_j>x_0$，两个列表没有重合。故这个集合恰有 $s+t$ 个元素。任何匹配每增加一对，都必须占用一个尚未用过的右侧整数，因此匹配至多有 $s+t$ 对。

我们已经给出了 $st$ 个不同正整数组成的集合 $A$，以及长度恰为 $2\max A$ 的开区间 $I$，其中任何匹配都不超过 $s+t$ 对。这正是定理~\ref{thm:product-upper} 要求的反向限制，于是
\[
\boxed{f(st)\leq s+t}.
\]

\begin{remark}\label{up:upper-only}
这里已经完成上界证明，不必再证明这个构造一定能配出 $s+t$ 对。一个保证值的上界，只要求给出无法超过该数的实例；要证明保证值恰好等于某个数，还必须另外证明任意集合、任意合适区间都至少能配出那么多对。两类证明中的量词不同，本章处理的是前者。我们把全部候选倍数逐一算清，是为了保证“至多”这个判断没有漏掉其他可以使用的整数。
\end{remark}

\section{完整数值实例：两行两列怎样实际算出来？}\label{up:example}

\begin{example}\label{up:two-by-two}
取 $s=t=2$，于是 $st=4$，并有
\[
D=\lcm(1,2,3,4)=12.
\]
这里 $q$ 只能取 $-1,1$，$r$ 只能取 $-1,0,1$，所以 $q+12r$ 的所有可能值为
\[
-13,-11,-1,1,11,13.
\]
其中大于 $4$ 的素因子只有 $11,13$，因此 $\mathcal P=\{11,13\}$。四个禁止表达式 $-i-12j$ 分别为 $-13,-25,-14,-26$；把它们分别按两个模数取余，得到
\[
\begin{array}{c|c|c}
\text{模数}&\text{禁止的不同余数}&\text{本例选择的允许余数}\\\hline
11&7,8,9&2\\
13&0,1,12&10
\end{array}
\]
选择 $M=101$，因为
\[
101\equiv2\pmod{11},\qquad101\equiv10\pmod{13},
\]
而且 $101>100$，也满足构造真正需要的条件
\[
101>2s+2tD=4+48=52.
\]
此时，按行号 $i$、列号 $j$ 排列的数表为
\[
\begin{array}{c|cc}
 &j=1&j=2\\\hline
i=1&\alpha_{1,1}=114&\alpha_{1,2}=126\\
i=2&\alpha_{2,1}=115&\alpha_{2,2}=127
\end{array}
\]
因此 $A=\{114,115,126,127\}$，最大值为 $127$。

下面实际求出 $x_0$，而不是只引用存在性。第一行要求 $x_0-1$ 同时被 $114,126$ 整除；二者的最小公倍数为 $2394$。第二行要求 $x_0-2$ 同时被 $115,127$ 整除；二者互素，所以最小公倍数为 $115\cdot127=14605$。于是方程组等价于
\[
x_0\equiv1\pmod{2394},\qquad
x_0\equiv2\pmod{14605}.
\]
令 $x_0=1+2394k$，就只需解 $2394k\equiv1\pmod{14605}$。带余除法给出
\begin{align*}
14605&=6\cdot2394+241,\\
2394&=9\cdot241+225,\\
241&=225+16,\\
225&=14\cdot16+1.
\end{align*}
从最后一行逐步代回，有
\begin{align*}
1&=225-14\cdot16\\
 &=15\cdot225-14\cdot241\\
 &=15\cdot2394-149\cdot241\\
 &=909\cdot2394-149\cdot14605.
\end{align*}
所以可以取 $k=909$，得到
\[
x_0=1+2394\cdot909=2176147.
\]
同时 $x_0=2+14605\cdot149$，故两个同余条件均可直接核对。四个原模数的要求也全部满足，因为 $114,126$ 都整除 $2394$，而 $115,127$ 都整除 $14605$。

现在设
\[
x=x_0-M=2176046,\qquad
I=(2176046,2176300).
\]
区间长度为 $2176300-2176046=254=2\cdot127$，正好符合定义。两个列表是
\[
L_1=2176146,\qquad L_2=2176145,
\]
以及
\[
R_1=2176260,\qquad R_2=2176272.
\]
它们与左端点的距离依次为 $100,99,214,226$，全都严格介于 $0$ 和 $254$ 之间，所以四点都在开区间内。每个格子对应的两个倍数可以直接写成乘法等式：
\[
\begin{array}{c|c|c}
\text{集合中的整数}&\text{区间内的第一个倍数}&\text{区间内的第二个倍数}\\\hline
114&114\cdot19089=2176146&114\cdot19090=2176260\\
126&126\cdot17271=2176146&126\cdot17272=2176272\\
115&115\cdot18923=2176145&115\cdot18924=2176260\\
127&127\cdot17135=2176145&127\cdot17136=2176272
\end{array}
\]
为了连“没有漏掉第三个倍数”也落实到数字，下表分别列出再往前、再往后的相邻倍数：
\[
\begin{array}{c|r|r}
\text{集合中的整数}&\text{前一个倍数}&\text{后一个倍数}\\\hline
114&2176032&2176374\\
126&2176020&2176398\\
115&2176030&2176375\\
127&2176018&2176399
\end{array}
\]
中间一列全部小于左端点 $2176046$，最后一列全部大于右端点 $2176300$。因此，区间中所有能被 $A$ 中某个元素整除的整数，完整列表就是
\[
\{2176145,2176146,2176260,2176272\}.
\]
这些数字虽大，但每一步只涉及整除、取余和整数加减乘除，均可以逐项复算。数值核验的作用是检查这个具体实例；任意 $s,t$ 的结论已经由前面的论证证明，不能用一个实例代替。

人工复核时，可以按三个层次进行。先检查余数：$M$ 是否避开两个坏素数对应的禁止余数，$x_0$ 是否分别比第一行的公倍数多 $1$、比第二行的公倍数多 $2$。再检查区间：长度必须是最大元素的两倍，而四个候选点必须严格位于端点之间。最后检查相邻倍数：既然每一行列出的两点相差恰好等于对应的整数，只要紧邻它们之外的两个倍数都在区间外，就已经排除了这个整数的全部其他倍数，无须逐个扫描区间中的每个整数。这三个层次分别对应一般证明中的同余、位置和完整性。
\end{example}

在两行两列时，$st=s+t=4$，所以这个实例主要展示构造怎样落地，还没有展示右侧数量比左侧严格少。当 $s=t=3$ 时，同一个一般构造就会把 $9$ 个不同正整数的区间内倍数限制在 $6$ 个点上；这句话直接应用已证明的定理，并不需要把这个更大的实例全部列出来。

\section{回看整条证明：每一步究竟解决什么？}\label{up:roadmap}

本章从一个可以检查的目标出发：让每个格子中的整数，恰好连接一个按行编号的倍数和一个按列编号的倍数。为此先令两个候选倍数之差等于该格子里的整数，再把“它们确实都是倍数”转化为一组同余方程。最小公倍数 $D$ 控制小素数的指数，参数 $M$ 避开有限多个大素数，二者合起来保证所有同余要求兼容。广义中国剩余定理于是给出共同的 $x_0$，最后用明确的端点不等式排除第三个倍数。

整个上界证明可以按以下逻辑链复核：
\[
\begin{gathered}
D\text{ 含有需要的小素数幂，}\quad M\text{ 避开坏素数}\\
\Longrightarrow\quad
\gcd(\alpha_{i,j},\alpha_{k,l})\mid(i-k)\\
\Longrightarrow\quad
\text{同一个 }x_0\text{ 满足全部余数要求}\\
\Longrightarrow\quad
I\text{ 内每个 }\alpha_{i,j}\text{ 的倍数恰为 }L_i,R_j\\
\Longrightarrow\quad
\text{全部可用右侧整数只有 }s+t\text{ 个}\\
\Longrightarrow\quad f(st)\leq s+t.
\end{gathered}
\]
前面的数论保证构造存在，后面的区间计算保证构造具有需要的性质。两部分缺少任何一项，都不能从一幅行列安排图直接得到上界。

还可以检查构造的先后依赖，确认它没有循环：先给定 $s,t$，就能计算 $D$；再由有限范围的 $q,r$ 确定坏素数，选择 $M$；有了 $M$ 才确定所有模数 $\alpha_{i,j}$，证明它们兼容后再选 $x_0$；最后才定义区间端点。坏素数的定义中没有未知的 $x_0$，选择 $M$ 时也没有提前使用尚未得到的同余解。因此，每一步都建立在已经确定的整数之上，这也是它能够写成具体数值例子的原因。

% ===== rounding =====
\chapter{上下界怎样合成精确答案}
\label{rnd:chapter}

\section{\texorpdfstring{先说明 $f(m)$ 为什么随 $m$ 增大而不减}{先说明 f(m) 为什么随 m 增大而不减}}

回顾定义：$f(m)$ 是对所有允许的集合和区间都成立的最大配对保证数。也就是说，无论怎样选取由 $m$ 个不同正整数组成的集合 $A$，也无论怎样移动长度为 $2\max A$ 的开区间，都至少能配成 $f(m)$ 对；这个统一保证不能再提高。一次匹配最多使用 $A$ 的全部 $m$ 个元素，所以始终有
\begin{equation}\label{rnd:trivial-upper}
f(m)\leq m.
\end{equation}
这里的保证涉及所有集合和所有区间。要证明 $f(m)\leq f(n)$，应当说明：任意一个 $n$ 元问题中，都能找到至少 $f(m)$ 对。不能只比较某两个恰好选中的集合。

\begin{lemma}[$f$ 的单调性]\label{rnd:monotonicity}
若 $1\leq m\leq n$，则 $f(m)\leq f(n)$。
\end{lemma}

\begin{proof}
任取一个含 $n$ 个不同正整数的集合 $A$，记 $a=\max A$，再任取实数 $x$。我们要在
\[
I=(x,x+2a)
\]
中找到至少 $f(m)$ 对。选取 $A$ 的任意一个 $m$ 元子集 $S$，记 $d=\max S$。因为 $S\subseteq A$，有 $d\leq a$。

如果只说“在 $S$ 上应用 $f(m)$ 的定义”，还没有交代区间长度：那个定义要求长度恰好为 $2d$，而原区间的长度为 $2a$。为此，截取同一个左端点的较短开区间
\[
J=(x,x+2d).
\]
由于 $d\leq a$，任何 $b\in J$ 都满足 $x<b<x+2d\leq x+2a$，所以 $J\subseteq I$。即使 $d=a$，两区间相等，也没有端点问题。

根据 $f(m)$ 的定义，集合 $S$ 与区间 $J$ 至少能配成 $f(m)$ 对。这些左点仍属于 $A$，右点仍位于 $I$，整除关系和双方各自互异的性质都没有改变，因而也是原问题的一组合法配对。因为最初的 $A$ 和 $x$ 完全任意，$n$ 元问题总能保证 $f(m)$ 对。因此 $f(n)\geq f(m)$。
\end{proof}

这个论证中，子集的最大元素可以变小；相应地，我们把区间也缩短。这样的配套处理，是把“在小集合中得到的保证”转移到大集合中所必需的。

\section{两个参数的上界怎样覆盖每个整数}

前面构造得到的定理~\ref{thm:product-upper} 告诉我们：对所有正整数 $s,t$，
\begin{equation}\label{rnd:product-bound}
f(st)\leq s+t.
\end{equation}
若 $m$ 恰好等于一个乘积，就可以直接应用；如果它的两个因子相差过大，则先把它增大到一个合适的乘积，再用引理~\ref{rnd:monotonicity}。例如 $14=2\times7$ 只能直接给出上界 $9$，而 $14\le16=4\times4$ 给出更好的上界 $8$。乘积应当选得既不小于 $m$，两个因子又尽量接近。

边界 $s=1$ 或 $t=1$ 不需要复杂构造。例如 $s=1$ 时，式~\eqref{rnd:trivial-upper} 直接给出 $f(st)=f(t)\leq t<1+t=s+t$。因此，下面出现因子 $1$ 时，式~\eqref{rnd:product-bound} 仍然有效。

还要单独认清 $m=1$。设 $A=\{a\}$，取
\[
q=\floor{x/a}+1.
\]
由向下取整的定义，$x/a<q\leq x/a+1$，乘以正数 $a$ 得
\[
x<qa\leq x+a<x+2a.
\]
所以任何要求的开区间中都至少有一个 $a$ 的倍数，可以配成一对；再由 $f(1)\leq1$ 得 $f(1)=1$。这里没有要求 $q$ 为正，因为区间中的倍数本来就允许是任意整数。

\begin{proposition}\label{rnd:ceiling-upper}
对每个正整数 $m$，都有 $f(m)\leq\ceil{2\sqrt m}$。
\end{proposition}

\begin{proof}
\textbf{情形一：$m$ 是完全平方数。}
写成 $m=k^2$，其中 $k\geq1$ 是整数。在式~\eqref{rnd:product-bound} 中取 $s=t=k$，得到
\[
f(m)=f(k^2)\leq2k=\ceil{2\sqrt m}.
\]
这也包括 $k=1$；此时的上界是 $2$，真正更强的上界 $1$ 会在最后取最小值时保留下来。

\textbf{情形二：$m$ 不是完全平方数，且 $m\leq k(k+1)$。}
令 $k=\floor{\sqrt m}$。因为 $m\geq1$ 且不是平方数，$k\geq1$，并且
\[
k^2<m<(k+1)^2.
\]
本情形假设 $k^2<m\leq k(k+1)$。由单调性及乘积上界，
\[
f(m)\leq f(k(k+1))\leq k+(k+1)=2k+1.
\]
要确认这正是要求的向上取整值，可以比较平方：
\[
(2k)^2<4m\leq4k(k+1)=(2k+1)^2-1<(2k+1)^2.
\]
各量非负，取平方根得到 $2k<2\sqrt m<2k+1$。因此 $\ceil{2\sqrt m}=2k+1$。

\textbf{情形三：$m$ 不是完全平方数，且 $m>k(k+1)$。}
仍取同一个 $k$。把 $m$ 增大到 $(k+1)^2$，得
\[
f(m)\leq f((k+1)^2)\leq2(k+1)=2k+2.
\]
由于 $m$ 是整数，严格不等式 $m>k(k+1)$ 给出 $m\geq k(k+1)+1$，所以
\[
4m\geq4k(k+1)+4=(2k+1)^2+3>(2k+1)^2.
\]
另一方面，$m<(k+1)^2$ 给出 $4m<(2k+2)^2$。取平方根可得
\[
2k+1<2\sqrt m<2k+2,
\]
从而 $\ceil{2\sqrt m}=2k+2$。三个情形包含全部正整数，命题得证。
\end{proof}

在这套分段中，$k(k+1)$ 属于第二种情形，紧接着的整数 $k(k+1)+1$ 属于第三种情形，两者没有遗漏或重叠。完全平方数则放在第一种情形中。也可以把第三段写到 $(k+1)^2$ 为止，但要同时调整第一段，以免重复讨论。

\begin{example}\label{rnd:small-upper-examples}
$m=6$ 时，$2^2<6=2\times3$，所以采用相邻因子 $2,3$，得到 $f(6)\leq5$。$m=7$ 时，$2\times3<7<3^2$，就不能再用六元集合提供七元问题的上界；应改用九元问题，得到 $f(7)\leq f(9)\leq6$。$m=10$ 时，$3^2<10\leq3\times4$，得到 $f(10)\leq f(12)\leq7$。

这些比较都使用了已证明的单调性。它们不要求为七或十硬凑一个乘积分解，也没有忽略“改变集合大小会连带改变区间长度”这一点；关于区间长度的处理，已经包含在引理~\ref{rnd:monotonicity} 的证明之中。
\end{example}

\section{为什么正好选择相等或相邻的两个因子}

上述选法并非偶然。我们已有的估计方式是：先选正整数 $s,t$，使 $st\geq m$，再推出
\[
f(m)\leq f(st)\leq s+t.
\]
因此，对这一类估计来说，目标是让 $s+t$ 尽量小。由平方非负，
\[
(s-t)^2\geq0
\quad\Longrightarrow\quad
(s+t)^2\geq4st\geq4m,
\]
所以任何这样的因子选择都满足 $s+t\geq2\sqrt m$。又因为因子和为整数，必有 $s+t\geq\ceil{2\sqrt m}$。前一节给出的选择恰好达到这个值，因此它在这类方法中已经最优。

也可以反过来，固定因子和，问乘积最多有多大。恒等式
\[
st=\frac{(s+t)^2-(s-t)^2}{4}
\]
说明：和固定后，两个因子越接近，差的平方就越小，乘积便越大。若和是偶数 $2k$，两个因子可以相等，最大乘积为 $k^2$；若和是奇数 $2k+1$，两个整数因子不能相等，差的绝对值至少为 $1$，最大乘积为
\[
\frac{(2k+1)^2-1}{4}=k(k+1).
\]
因此，当允许的因子和依次是 $2k,2k+1,2k+2$ 时，能够覆盖的最大乘积恰好依次是 $k^2,k(k+1),(k+1)^2$。这正是分段证明使用的三个边界。

需要分清这段论证的作用：它证明了“利用乘积上界选择因子”的方式已经不能给出更小的数，却没有单独证明真正的 $f(m)$ 不能更小。关于 $f(m)$ 的下界，仍须依靠前面的匹配与邻域论证。一个上界方法达到了它自身的极限，不等于对象本身的精确值已经确定。

\section{整数取整与最小值各自承担什么作用}

向上取整 $\ceil y$ 表示不小于 $y$ 的最小整数。因此，如果一个整数 $r$ 满足 $r\geq y$，就能加强为 $r\geq\ceil y$。反过来，如果整数 $r\leq y$，对应的加强是 $r\leq\floor y$；不等号方向不能混淆。本题下界证明先得到一个实数下界，再利用匹配数为整数而向上取整；刚才的上界则由具体构造直接给出所需的整数。

对整数 $m$ 和实数 $y$，确有恒等式
\begin{equation}\label{rnd:min-ceiling}
\ceil{\min(m,y)}=\min(m,\ceil y).
\end{equation}
若 $y\leq m$，因为 $m$ 本身是一个不小于 $y$ 的整数，所以 $\ceil y\leq m$，两边都等于 $\ceil y$；若 $y>m$，两边都等于 $m$。这说明在这里可以把取整写到 $y$ 上，同时保留整数 $m$。并非看到取整符号就能对里面的所有运算随意拆开。

尤其要区分
\[
\ceil{2\sqrt m}
\quad\text{与}\quad
2\ceil{\sqrt m}.
\]
以 $m=5$ 为例，因为 $4<2\sqrt5<5$，前者为 $5$；因为 $2<\sqrt5<3$，后者为 $2\times3=6$。把倍数 $2$ 移到取整符号外，会失去一个单位的精度。这个单位恰好会影响精确答案。

\begin{theorem}[精确答案]\label{rnd:exact-answer}
对每个正整数 $m$，
\[
\boxed{f(m)=\min\bigl(m,\ceil{2\sqrt m}\bigr).}
\]
\end{theorem}

\begin{proof}
下界部分已经对任意由 $m$ 个不同正整数组成的集合，以及任意允许的开区间证明
\[
f(m)\geq\min\bigl(m,\ceil{2\sqrt m}\bigr).
\]
另一方面，式~\eqref{rnd:trivial-upper} 给出 $f(m)\leq m$，命题~\ref{rnd:ceiling-upper} 给出 $f(m)\leq\ceil{2\sqrt m}$。同时满足这两个上界，正是
\[
f(m)\leq\min\bigl(m,\ceil{2\sqrt m}\bigr).
\]
同一个整数被夹在两个相同的数之间，因而只能等于它们。
\end{proof}

当 $m\geq4$ 时，$\sqrt m\geq2$，乘以 $\sqrt m$ 得 $m\geq2\sqrt m$。由于 $m$ 是整数，进一步有 $m\geq\ceil{2\sqrt m}$，所以此时答案可以简写成 $f(m)=\ceil{2\sqrt m}$。而当 $m=1,2,3$ 时，$\ceil{2\sqrt m}$ 分别为 $2,3,4$，均超过可供配对的左点总数，必须取最小值，得到 $f(m)=m$。

\section{从 1 到 16 的精确数值}

表中的平方根取整不依靠小数近似：$\ceil{2\sqrt m}$ 就是满足 $r^2\geq4m$ 的最小非负整数 $r$。因此用整数平方比较就可以逐项检查。比如 $m=12$ 时，$6^2=36<48\leq49=7^2$，故取整值为 $7$；$m=13$ 时，$7^2=49<52\leq64=8^2$，取整值升为 $8$。

\begin{center}
\renewcommand{\arraystretch}{1.2}
\begin{tabular}{c|c|c@{\qquad}c|c|c}
\hline
$m$ & $\ceil{2\sqrt m}$ & $f(m)$ & $m$ & $\ceil{2\sqrt m}$ & $f(m)$\\
\hline
1 & 2 & 1 & 9 & 6 & 6\\
2 & 3 & 2 & 10 & 7 & 7\\
3 & 4 & 3 & 11 & 7 & 7\\
4 & 4 & 4 & 12 & 7 & 7\\
5 & 5 & 5 & 13 & 8 & 8\\
6 & 5 & 5 & 14 & 8 & 8\\
7 & 6 & 6 & 15 & 8 & 8\\
8 & 6 & 6 & 16 & 8 & 8\\
\hline
\end{tabular}
\end{center}

其中 $f(6)=5$ 的含义是：任意六个不同正整数和任意规定长度的开区间，总能找到五对互不冲突的配对；同时确实存在某些选择，使六对不可能同时实现。下界证明负责“无论怎样选都能做到”，上界构造负责“有一种选择不能更多”。两部分合在一起，才确定了这个最优保证值。

% ===== examples =====
\chapter{把抽象证明落到具体计算上}
\section{六个数确实可能只能配成五对}
主定理给出 $f(6)=\ceil{2\sqrt6}=5$。下面写出一个完全具体的情形，展示为什么不能把保证值提高为六对。

设
\[
N=69898816,\qquad A=\{51,52,57,58,63,64\},
\]
并取
\[
I=(N-86,N+42)=(69898730,69898858).
\]
最大数是 $64$，区间长度为 $86+42=128=2\times64$，符合题目的要求。虽然 $N$ 较大，接下来的结构只需要记住五个小的偏移量。

先直接核对两个整数的整除关系：
\[
\begin{aligned}
N-1&=51\times1370565=57\times1226295=63\times1109505,\\
N&=52\times1344208=58\times1205152=64\times1092169.
\end{aligned}
\]
于是每个 $a$ 的两个连续倍数如下。
\begin{center}
\begin{tabular}{ccc}
\toprule
$a$ & 区间内的较小倍数 & 区间内的较大倍数\\
\midrule
$51$ & $N-52$ & $N-1$\\
$52$ & $N-52$ & $N$\\
$57$ & $N-58$ & $N-1$\\
$58$ & $N-58$ & $N$\\
$63$ & $N-64$ & $N-1$\\
$64$ & $N-64$ & $N$\\
\bottomrule
\end{tabular}
\end{center}
每行两个数的差恰好是左边的 $a$，所以它们是连续的两个倍数。它们都严格位于 $I$ 内：最小的是 $N-64>N-86$，最大的是 $N<N+42$。

还必须证明没有第三个倍数。较小倍数最多为 $N-52$，而每个 $a\ge51$，所以它前面的那个倍数至多为
\[
(N-52)-51=N-103<N-86.
\]
较大倍数至少为 $N-1$，因此它后面的那个倍数至少为
\[
(N-1)+51=N+50>N+42.
\]
更远的倍数自然在这两个外侧位置之外。因此上表完整列出了区间内的全部候选。

六个数合起来，只能用到五个整数
\[
N-64,\quad N-58,\quad N-52,\quad N-1,\quad N.
\]
同一个整数不能被用两次，所以最多配成五对。另一方面，下列五对确实可同时使用：
\[
(51,N-1),\quad(58,N),\quad(52,N-52),\quad
(57,N-58),\quad(63,N-64).
\]
第一列互异，第二列也互异，每一对的整除关系都能由上表核对。这个具体情形的最大配对数因此\textbf{恰好是五}。

这里用程序辅助找到了一个便于列出完整表格的例子；它成立的依据是上面的逐项整除与边界计算，而不是“程序说可以”。上界章节已经给出了对任意正整数参数 $s,t$ 都适用的构造证明。

\section{区间长度的系数为什么重要}
为了看清题目的尺度，暂把区间长度改成 $cL$，并把相应的统一保证值记作 $f_c(m)$。原问题就是 $c=2$。

\subsection{\texorpdfstring{当 $1<c<2$ 时，只能统一保证一对}{当 1<c<2 时，只能统一保证一对}}
任取 $A,x$。令 $a$ 是 $A$ 中任意一个数，并取
\[
b=a\left(\floor{x/a}+1\right).
\]
由取整定义有 $x<b\le x+a$。又因为 $a\le L<cL$，所以 $b<x+cL$。这样至少能配成一对。

现在构造一个只能配成一对的例子。取充分大的正整数 $M$，令
\[
A=\{M+1,\ldots,M+m\},\qquad
P=\prod_{i=1}^m(M+i),
\]
并选区间
\[
I_c=(P-(M+1),\ P-(M+1)+c(M+m)).
\]
我们要求
\[
c(M+m)<2(M+1),
\quad\text{即}\quad (2-c)M>cm-2.
\]
由于 $2-c>0$，可以选一个整数 $M$ 满足这条严格不等式。$P$ 位于区间内部：左端点小于 $P$，而 $c(M+m)>M+1$ 保证右端点大于 $P$。

每个 $a\in A$ 都整除 $P$。它前面的倍数 $P-a$ 不大于左端点 $P-(M+1)$；它后面的倍数 $P+a$ 不小于 $P+(M+1)$，又严格大于区间右端点。因此区间里唯一能被 $A$ 中任何一个数整除的整数就是 $P$。所有数都只能争这一个位置，最多形成一对。结合前面的保证，得到 $f_c(m)=1$。

\subsection{\texorpdfstring{当 $2\le c<3$ 时，本书的答案仍然成立}{当 2≤c<3 时，本书的答案仍然成立}}
较长区间 $(x,x+cL)$ 包含 $(x,x+2L)$，所以原来的下界仍然成立。

上界构造也可以沿用，但要把 $M$ 选得更大。令 $C=s+tD$，则最大数为 $M+C$。要求
\[
(3-c)M>cC,
\quad\text{即}\quad M>\frac{c}{3-c}(s+tD).
\]
这里 $3-c>0$，所以这仍只是一个可满足的“足够大”条件。由于 $c\ge2$，它也保证 $M>2(s+tD)$，原来两类候选点都在短区间内部，因此也在长区间内部。

位于两类候选点之前的相邻倍数仍在左端点外；第三个倍数满足
\[
(x_0-i)+2\alpha_{i,j}>x+3M
>x+c(M+s+tD).
\]
所以扩大区间后依然没有新的可用倍数。坏素数避让、最大公约数兼容与余数方程都不受影响，因为那些步骤允许在同一个余数类中任取充分大的 $M$。于是仍有 $f_c(st)\le s+t$；上一章的单调性和取整推导原样适用，得到
\[
f_c(m)=\min\{m,\ceil{2\sqrt m}\},\qquad 2\le c<3.
\]
当 $c\ge3$ 时，刚才使第三个倍数留在区间外的条件不再成立，这一构造不能直接给出相同结论。本讲义只证明上述两个明确的范围。

\section{几个容易混淆的地方}
\subsection{有很多邻点，不等于已经选好了匹配}
下界中给每个 $a$ 指派的两个倍数，允许不同的 $a$ 共用某个坐标。例如 $(6,12)$ 与 $(8,12)$ 是两个不同的有序对，它们却共用 $12$。单射保证的是\emph{整个有序对}不重复。它先产生邻点数量的不等式，再由配对准则得到互不冲突的匹配，两个步骤缺一不可。

\subsection{为什么必须检查所有子集}
一整组数的候选倍数很多，仍可能存在一小组数挤在很少的位置上。亏损公式把这件事准确写成
\[
\max_{S\subseteq A}(|S|-|\Gamma(S)|).
\]
最大亏损可能来自真子集，不能只计算 $S=A$。

\subsection{中间估计有例外，不表示最终定理失效}
取 $L=10,x=0$，并令 $S=\{10\}$。开区间 $(0,20)$ 内只有一个 $10$ 的倍数，即 $10$。于是 $|\Gamma(S)|=1$，而 $2\sqrt{|S|}=2$。这说明不能无条件地对原图套用邻域估计。

正文的实际处理是先固定配对 $(10,10)$，删去这两个顶点，再在剩余图上证明邻域估计。删去特殊点和加回这条边，都已经包含在下界证明中。因此这个例子恰好说明分类处理为什么必要。

\subsection{检验有限样本不能替代一般证明}
程序可以检查几万个具体整数，帮助发现端点、取整或分类错误。但是 $m,A,x$ 有无限多种可能，有限样本本身不能覆盖所有输入。正文使用任意的参数作推导，这才是一般证明；计算核验是另一个层面的检查。

\section{读完后可以自己完成的练习}
\begin{enumerate}
\item 对 $A=\{2,3,4\}$、$I=(0,8)$，列出各个数的候选倍数，并给出三对互不冲突的配对。
\item 已知 $p,q$ 是非负整数且 $pq\ge13$，证明 $p+q\ge8$。
\item 取 $L=10,b_0=10$。分别求 $a=2$ 和 $a=4$ 在修正后的三类映射下的像，解释朴素选点为什么碰撞。
\item 对非空集合大小 $n=10$，若所有子集满足邻域估计，本书的准则保证至少多少对？
\item 解余数方程 $z\equiv1\pmod6$、$z\equiv3\pmod8$；再说明把第二条改为 $z\equiv2\pmod8$ 后为什么无解。
\item 分别计算 $f(10),f(12),f(13)$，并指出各自使用哪个乘积上界。
\item 设 $q$ 为非零整数，$s,t$ 为正整数。解释为什么 $q^2<st$ 且 $D=\lcm(1,\ldots,st)$ 能推出 $q^2\mid D$，但 $|q|\mid D$ 单独不能推出 $q^2\mid D$。
\item 为什么得到“匹配大小至少 $2\sqrt m$”之后，应向上取整，而不能向下取整？
\end{enumerate}

\section{练习的完整答案}
\begin{enumerate}
\item 候选分别为 $\{2,4,6\}$、$\{3,6\}$、$\{4\}$。选择 $(2,2),(3,3),(4,4)$，每一列都互异，且每行满足整除要求。
\item 若 $p+q\le7$，由 $(p-q)^2\ge0$ 得 $pq\le49/4=12.25$，与 $pq\ge13$ 矛盾。因此 $p+q\ge8$。
\item $2\mid10$，且 $4<10$、$4\nmid10$，所以 $a=2$ 属于第二类，像为 $(6,14)$。$4\nmid10$，所以 $a=4$ 属于第一类，像为 $(8,12)$。若把 $a=2$ 也简单地放在 $10\pm2$，就会得到同样的 $(8,12)$，造成碰撞。
\item 保证值是 $\min(10,\ceil{2\sqrt{10}})=7$，因为 $6^2=36<40\le49=7^2$。
\item 所有形如 $z=1+6k$ 的整数中，要求 $1+6k\equiv3\pmod8$，即 $3k\equiv1\pmod4$，故 $k\equiv3\pmod4$。因此 $z\equiv19\pmod{24}$，最小非负解为 $19$。第二组余数方程要求一个数同时为奇数和偶数，无解；等价地，$\gcd(6,8)=2$ 不整除余数差 $1-2=-1$。
\item $f(10)=7$、$f(12)=7$、$f(13)=8$。前两者用 $10\le3\times4$ 和 $12=3\times4$，于是上界为 $3+4=7$；最后用 $13\le4\times4$，上界为 $8$。下界章分别保证相同数量。整数平方比较为 $36<40\le49$、$36<48\le49$、$49<52\le64$。
\item 因为 $q\ne0$，$q^2$ 是 $1,\ldots,st$ 中的一个整数，而 $D$ 被定义为这整个列表的公共倍数，所以 $q^2\mid D$。如果只知道 $|q|\mid D$，结论不成立，例如 $2\mid6$ 但 $4\nmid6$。
\item 匹配大小是整数。若一个整数至少为 $6.3$，它最少只能是 $7$，不可能是 $6$。一般地，整数 $r\ge y$ 等价于 $r\ge\ceil y$；这正是向上取整的定义。
\end{enumerate}

\backmatter
\chapter{来源与核查范围}
原论文的第2节陈述主定理与匹配公式，第3节给出上界构造，第4节给出含端点分类的下界证明。本讲义的“任意情形能做到”和“特殊情形不能更多”分别对应原文这两部分。配对公式的交错路证明、剩余定理的初等铺垫、具体数值表及练习，是为阅读补充的说明。

数学核查重点包括：开区间不能取到端点；中点是最大数的倍数时必须另行处理；三类映射的三种跨类碰撞都要排除；广义剩余定理必须检查最大公约数兼容条件；单调性证明必须处理集合最大值变化导致的区间长度变化；小的 $m$ 和取整边界必须单列。

随包的 \texttt{check\_examples.py} 用有理数与整数做有限范围核验，包含一般情形的选点、端点情形的三类映射、具体上界实例和取整值。其通过仅表示这些有限样本没有发现错误。对所有输入成立的依据是正文的证明。详细的实际运行记录见随包 \texttt{PROOF\_CHECKS.md}。

\begin{thebibliography}{9}
\addcontentsline{toc}{chapter}{参考文献}
\bibitem{paper} Wouter van Doorn, Yanyang Li, Quanyu Tang.
\emph{Optimal bounds for an Erd\H{o}s problem on matching integers to distinct multiples}.
arXiv:2603.28636v1, 2026-03-30.
\url{https://arxiv.org/abs/2603.28636v1}.
\bibitem{translation} 同文的完整中文译本与原文底本，Erd\H{o}s650 中文翻译工程。
\url{https://github.com/liyiliuxingyu/erdos650-chinese-translation}.
\end{thebibliography}
\end{document}
